\documentclass[12pt,a4paper,twoside,openany]{report}

\usepackage{fontspec}
\usepackage{polyglossia}
\setmainlanguage{slovak}
\setotherlanguage{english}
\setmainfont{Times New Roman}
\setsansfont{Arial}
\setmonofont{Consolas}

% Odporúčané nastavenie podľa IS 1/2024 (čl. 7):
% vnútorný okraj 3,5 cm; vonkajší okraj 2 cm; horný a dolný 2,5 cm; päta 1,25 cm.
% Pri obojstrannej tlači sa používajú zrkadlové okraje.
\usepackage[
    a4paper,
    inner=35mm,
    outer=20mm,
    top=25mm,
    bottom=25mm,
    footskip=12.5mm
]{geometry}

\usepackage{setspace}
\onehalfspacing
\usepackage{indentfirst}
\setlength{\parindent}{1.25cm}
\setlength{\parskip}{0pt}
\usepackage[nopatch=footnote]{microtype}
\usepackage{graphicx}
\usepackage{float}
\usepackage{enumitem}
\usepackage{needspace}
\usepackage[section]{placeins}
\usepackage{amsmath,amssymb}
\usepackage{csquotes}
\usepackage{tikz}
\usetikzlibrary{arrows.meta,positioning,fit,calc,shapes.geometric}
\usetikzlibrary{backgrounds,decorations.pathreplacing}
\usepackage[hidelinks,pdfpagelayout=SinglePage]{hyperref}
\usepackage[indentafter]{titlesec}
\usepackage[
    backend=biber,
    style=apa,
    sorting=nyt
]{biblatex}
\addbibresource{references.bib}

% Formátovanie nadpisov podľa odporúčania v smernici:
% 1. úroveň 16b/Bold, 2. úroveň 14b/Bold, 3. úroveň 14b/Italic.
\titleformat{\chapter}
  {\bfseries\fontsize{16pt}{19pt}\selectfont}
  {\thechapter}{1em}{}
\titleformat{\section}
  {\bfseries\fontsize{14pt}{17pt}\selectfont}
  {\thesection}{1em}{}
\titleformat{\subsection}
  {\itshape\fontsize{14pt}{17pt}\selectfont}
  {\thesubsection}{1em}{}
\titlespacing*{\chapter}{0pt}{0pt}{18pt}
\titlespacing{\section}{0pt}{18pt}{8pt}
\titlespacing{\subsection}{0pt}{14pt}{6pt}

% Jednotné nastavenie zoznamov v celej práci.
\setlist[itemize]{leftmargin=1.4cm,itemsep=4pt,topsep=5pt,parsep=0pt,partopsep=0pt}
\setlist[enumerate]{leftmargin=1.5cm,itemsep=4pt,topsep=5pt,parsep=0pt,partopsep=0pt}

% Kompaktný blok na vysvetlenie symbolov a skratiek po vzore "kde:".
\newenvironment{symboltable}
  {\begin{list}{}%
    {\setlength{\leftmargin}{1.1cm}%
     \setlength{\itemindent}{0pt}%
     \setlength{\labelwidth}{0pt}%
     \setlength{\labelsep}{0pt}%
     \setlength{\itemsep}{1.5pt}%
     \setlength{\topsep}{3pt}%
     \setlength{\parsep}{0pt}%
     \setlength{\partopsep}{0pt}}}
  {\end{list}}
\newcommand{\symbolitem}[2]{%
  \item \textendash\enspace #1\ #2
}

% Obsah v štýle vzoru (Príloha 4): "O B S A H" + "str." vpravo.
\makeatletter
\renewcommand{\contentsname}{O B S A H}
\renewcommand{\listfigurename}{Zoznam ilustrácií}
\renewcommand{\listtablename}{Zoznam tabuliek}
\renewcommand{\tableofcontents}{%
  \chapter*{\contentsname\hfill str.}%
  \@mkboth{\MakeUppercase\contentsname}{\MakeUppercase\contentsname}%
  \@starttoc{toc}%
}
% Ulozenie povodneho stylu "plain", aby sa dal obnovit od casti Uvod.
\let\ps@origplain\ps@plain
\makeatother

\newcommand{\UniversityName}{Ekonomická univerzita v Bratislave}
\newcommand{\FacultyName}{Fakulta hospodárskej informatiky}
\newcommand{\DepartmentName}{Katedra operačného výskumu a ekonometrie}

% Typ práce podľa zadania
\newcommand{\WorkTypeSK}{Diplomová práca}

\newcommand{\ThesisTitleSK}{Využitie umelej inteligencie pri riešení optimalizačných problémov}
\newcommand{\ThesisTitleEN}{Using artificial intelligence to solve optimization problems}
\newcommand{\ThesisSubtitleSK}{}

\newcommand{\AuthorName}{Bc. Oleksii Roiko}
\newcommand{\SupervisorName}{doc. Ing. Zuzana Čičková, PhD.}

% Prevezmi údaje z AIS2 a uprav podľa potreby.
\newcommand{\StudyProgramName}{Data science v ekonómii}
\newcommand{\StudyFieldName}{Ekonómia a manažment}
\newcommand{\EvidenceNumber}{103003/I/2026/36159551160335108}

\newcommand{\PlaceName}{Bratislava}
\newcommand{\SubmissionYear}{2026}
\newcommand{\FacultyAbbrev}{FHI EU}
\newcommand{\PageCount}{doplniť počet strán, napr. 70 s.}

% Abstrakt (podľa vzoru z Prílohy 3 v IS 1/2024)
\newcommand{\AbstractGoalSK}{preskúmať možnosti využitia moderných metód umelej inteligencie pri riešení symetrického problému obchodného cestujúceho a overiť ich praktickú použiteľnosť}
\newcommand{\AbstractDataSK}{synteticky generovaných inštancií a z referenčných súborov úloh}
\newcommand{\AbstractMethodSK}{modelovania hrán v grafoch pomocou neurónových sietí, vrátane porovnania dekódovacích postupov a lokálnej optimalizácie trasy}
\newcommand{\AbstractResultSK}{systematické porovnanie prístupov, ich silných a slabých stránok a podmienok, v ktorých sú použiteľné}
\newcommand{\AbstractValueSK}{prepojení celého experimentálneho cyklu od generovania údajov cez tréning po hodnotenie modelov a v zhrnutí odporúčaní pre ďalší výskum}
\newcommand{\AbstractGoalEN}{to examine the possibilities of applying modern artificial intelligence methods to the symmetric Traveling Salesman Problem and to verify their practical usability}
\newcommand{\AbstractDataEN}{synthetically generated instances and reference task sets}
\newcommand{\AbstractMethodEN}{edge modeling in graphs using neural networks, including a comparison of decoding procedures and local tour optimization}
\newcommand{\AbstractResultEN}{a systematic comparison of the approaches, their strengths and weaknesses, and the conditions under which they are applicable}
\newcommand{\AbstractValueEN}{linking the full experimental cycle from data generation through training to model evaluation, and summarizing recommendations for further research}
\newcommand{\KeywordsSK}{symetrický problém obchodného cestujúceho, neurónové siete, grafové modelovanie, strojové učenie, optimalizačné úlohy, lokálna optimalizácia trasy, referenčné súbory úloh}
\newcommand{\KeywordsEN}{symmetric Traveling Salesman Problem, neural networks, graph modeling, machine learning, optimization tasks, local tour optimization, reference task sets}

% Zadanie práce (z AIS2)
\newcommand{\AssignmentGoalSK}{Cieľom práce je preskúmať možnosti využitia neurónových sietí pri riešení symetrického problému obchodného cestujúceho a overiť ich praktickú použiteľnosť 
na syntetických a štandardných testovacích inštanciách. Práca sa zároveň zameriava na formuláciu úlohy v podobe vhodnej pre dátovo riadené modelovanie založené na predikcii hrán v úplnom grafe.

Na dosiahnutie uvedeného cieľa boli stanovené tieto úlohy:
\begin{itemize}
\item preskúmať teoretické východiská symetrického problému obchodného cestujúceho,
\item analyzovať existujúce prístupy k riešeniu tejto úlohy so zameraním na neurónové siete,
\item navrhnúť a realizovať experimentálny postup od prípravy syntetických údajov a štandardných testovacích inštancií cez trénovanie modelov až po hodnotenie a analýzu dosiahnutých výsledkov.
\end{itemize}
}


\begin{document}
\raggedbottom

% Číslovanie začína od titulného listu, ale zobrazí sa až od časti Úvod.
\pagenumbering{arabic}
\pagestyle{empty}
% Pred castou Uvod pouzijeme jednostrannu geometriu (bez zrkadlenia okrajov).
\newgeometry{
  asymmetric,
  left=35mm,
  right=20mm,
  top=25mm,
  bottom=25mm,
  footskip=12.5mm
}
% V uvodnej casti skryjeme cisla aj na stranach, kde chapter* vynucuje "plain".
\makeatletter
\let\ps@plain\ps@empty
\makeatother

\begin{titlepage}
\begingroup
\singlespacing
\centering

\vspace*{8mm}
{\Large\bfseries \MakeUppercase{\UniversityName}\par}
\vspace{2mm}
{\Large\bfseries \MakeUppercase{\FacultyName}\par}

\vspace{18mm}
\begin{flushleft}
Evidenčné číslo: \EvidenceNumber
\end{flushleft}

\vfill

{\LARGE\bfseries \ThesisTitleSK\par}
\vspace{8mm}
\ifx\ThesisSubtitleSK\empty
\else
{\Large \ThesisSubtitleSK\par}
\vspace{8mm}
\fi
\vspace{1mm}
{\large \WorkTypeSK\par}

\vfill
\noindent
\begin{minipage}[t]{0.48\textwidth}
\raggedright
\SubmissionYear
\end{minipage}
\begin{minipage}[t]{0.48\textwidth}
\raggedleft
\AuthorName
\end{minipage}

\endgroup
\end{titlepage}

\begin{titlepage}
\begingroup
\singlespacing
\centering

\vspace*{8mm}
{\Large\bfseries \MakeUppercase{\UniversityName}\par}
\vspace{2mm}
{\Large\bfseries \MakeUppercase{\FacultyName}\par}

\vfill

{\LARGE\bfseries \ThesisTitleSK\par}
\vspace{8mm}
\ifx\ThesisSubtitleSK\empty
\else
{\Large \ThesisSubtitleSK\par}
\vspace{8mm}
\fi
{\large \WorkTypeSK\par}

\vspace{10mm}

\begin{flushleft}
\textbf{Študijný program:} \StudyProgramName\par
\vspace{2mm}
\textbf{Študijný odbor:} \StudyFieldName\par
\vspace{2mm}
\textbf{Školiace pracovisko:} \DepartmentName\par
\vspace{2mm}
\textbf{Vedúca záverečnej práce:} \SupervisorName\par
\end{flushleft}

\vfill
\noindent
\begin{minipage}[t]{0.58\textwidth}
\raggedright
\PlaceName\ \SubmissionYear
\end{minipage}
\begin{minipage}[t]{0.40\textwidth}
\raggedleft
\AuthorName
\end{minipage}

\endgroup
\end{titlepage}


% Voliteľné strany:
\chapter*{Poďakovanie}
\addcontentsline{toc}{chapter}{Poďakovanie}

\noindent
Ďakujem vedúcej záverečnej práce doc. Ing. Zuzane Čičkovej, PhD., za odborné
vedenie, cenné pripomienky a podporu počas vypracovania tejto práce.
Zároveň ďakujem pedagógom Ekonomickej univerzity v Bratislave za vedomosti
a skúsenosti, ktoré mi počas štúdia odovzdali.

\clearpage

% \includepdf[pages=-]{frontmatter/zadanie.pdf}

\phantomsection
\addcontentsline{toc}{chapter}{Abstrakt}
\noindent{\bfseries\fontsize{16pt}{19pt}\selectfont ABSTRAKT\par}

\vspace{7mm}
\noindent
\textbf{\MakeUppercase{\AuthorName}:} \ThesisTitleSK. -- \UniversityName.
\FacultyName; \DepartmentName. -- Vedúci/Vedúca záverečnej práce:
\SupervisorName. -- \PlaceName: \FacultyAbbrev, \SubmissionYear, \PageCount.

\vspace{6mm}
\noindent
Cieľom záverečnej práce bolo \AbstractGoalSK.
Údaje sme získali z \AbstractDataSK, na ich spracovanie sme použili metódu
\AbstractMethodSK. Výsledkom riešenia práce je \AbstractResultSK.
Pridaná hodnota práce je v \AbstractValueSK.

\vspace{4mm}
\noindent
\textbf{Kľúčové slová:} \KeywordsSK

\clearpage

\phantomsection
\addcontentsline{toc}{chapter}{Abstract}
\noindent{\bfseries\fontsize{16pt}{19pt}\selectfont ABSTRACT\par}

\vspace{7mm}
\noindent
\textbf{\MakeUppercase{\AuthorName}:} \ThesisTitleEN. -- \UniversityName.
\FacultyName; \DepartmentName. -- Thesis supervisor:
\SupervisorName. -- \PlaceName: \FacultyAbbrev, \SubmissionYear, \PageCount.

\vspace{6mm}
\noindent
The objective of the thesis was \AbstractGoalEN.
Data were obtained from \AbstractDataEN, and processed using
\AbstractMethodEN. The main result of the thesis is \AbstractResultEN.
The added value of the thesis lies in \AbstractValueEN.

\vspace{3mm}
\noindent
\textbf{Keywords:} \KeywordsEN

\clearpage


\tableofcontents
\clearpage

% Voliteľné zoznamy podľa IS 1/2024.
% Zoznam ilustrácií je vhodné zaradiť už teraz, keďže práca už obsahuje obrázky.
\phantomsection
\addcontentsline{toc}{chapter}{\listfigurename}
\listoffigures
\clearpage

% Zoznam tabuliek zapneme po doplnení prvej tabuľky do práce.
% \phantomsection
% \addcontentsline{toc}{chapter}{\listtablename}
% \listoftables
% \clearpage

% Od casti Uvod zapneme geometriu so zrkadlovymi okrajmi.
\newgeometry{
  inner=35mm,
  outer=20mm,
  top=25mm,
  bottom=25mm,
  footskip=12.5mm
}
% Od casti Uvod obnovime standardny "plain" so zobrazenim cisel stran.
\makeatletter
\let\ps@plain\ps@origplain
\makeatother
\pagestyle{plain}

\chapter{Úvod}

\indent V súčasnosti zohrávajú úlohy obchodného cestujúceho (TSP) významnú úlohu v oblasti optimalizácie. S podobnými problémami sa stretávame v mnohých praktických situáciách a existuje veľké množstvo prístupov, ktoré umožňujú riešiť TSP aj jeho rôzne varianty. V praxi sa často používajú najmä heuristické a metaheuristické metódy, ktoré spravidla poskytujú suboptimálne riešenia.

V tejto práci sme sa na danú problematiku zamerali z pohľadu moderných metód umelej inteligencie. Motiváciou bolo overiť, či je možné využiť neurónové siete aj pri riešení TSP tak, aby boli výsledky použiteľné nielen z teoretického hľadiska, ale aj v praktických úlohách. Vychádzali sme z predpokladu, že aj napriek prirodzeným obmedzeniam môžu tieto metódy priniesť použiteľné riešenia a nové poznatky.

V rámci práce sme realizovali kompletný postup od generovania údajov, cez prípravu datasetov a trénovanie modelov, až po ich vyhodnotenie na referenčných benchmarkových dátach TSPLIB. Týmto postupom sme overili, že zvolený prístup dokázal riešiť úlohy TSP s merateľnou presnosťou a umožnil systematicky porovnať silné a slabé stránky jednotlivých modelových variantov.

Hlavnou motiváciou práce bolo tiež to, že umelá inteligencia sa dnes intenzívne zavádza do rôznych oblastí praxe. Preto sme považovali za dôležité overiť jej potenciál aj v optimalizačných úlohách tohto typu. Predložená práca tak môže slúžiť ako základ pre ďalší výskum a zároveň ako praktický rámec pre aplikáciu neurónových sietí pri riešení problémov TSP.

Cieľom práce bolo analyzovať možnosti riešenia optimalizačných úloh rôznymi prístupmi umelej inteligencie so zameraním na symetrický problém obchodného cestujúceho. Predmetom skúmania boli modely neurónových sietí určené na predikciu hrán v grafe, trénované na synteticky generovaných dátach a hodnotené na referenčných inštanciách TSPLIB. V práci sme sa sústredili na porovnanie kvality riešení pomocou optimalizačnej odchýlky voči referenčným trasám.

Prácu sme rozčlenili do viacerých logicky previazaných častí. V druhej kapitole sme predstavili teoretické východiská problematiky TSP, základné súvislosti umelej inteligencie a analýzu súčasného stavu riešenej problematiky. V tretej kapitole sme formulovali cieľ práce a podrobne opísali použitú metodiku, postup realizácie práce a nástroje využité pri experimentoch. Vo štvrtej kapitole sme uviedli dosiahnuté výsledky, ich porovnanie a diskusiu. V piatej kapitole sme zhrnuli hlavné závery práce a navrhli odporúčania pre ďalší výskum.

\chapter{Súčasný stav riešenej problematiky doma a v zahraničí}

V tejto kapitole sa zameriavame na teoretické a metodologické východiská symetrického problému obchodného cestujúceho a na hlavné smery jeho súčasného riešenia. Pozornosť venujeme najmä prechodu od klasických optimalizačných prístupov k dátovo riadeným metódam, ktoré využívajú neurónové siete a grafové reprezentácie. Keďže práca je orientovaná na modelovanie nad hranami úplného grafu, štruktúra kapitoly je prispôsobená práve tým aspektom problému, ktoré sú relevantné pre edge-based formuláciu, syntetické tréningové dáta, benchmarkové hodnotenie a interpretáciu výsledkov.

\section{Teoretické východiská problému obchodného cestujúceho}

Problém obchodného cestujúceho (Traveling Salesman Problem, TSP) patrí medzi najznámejšie úlohy kombinatorickej optimalizácie. V základnej formulácii je daná množina miest a náklady presunu medzi každou dvojicou miest, pričom cieľom je určiť takú uzavretú trasu, ktorá navštívi každé mesto práve raz a minimalizuje celkové náklady. Napriek relatívne jednoduchej formulácii ide o úlohu s vysokou výpočtovou náročnosťou, ktorá sa dlhodobo používa ako referenčný problém pri porovnávaní presných algoritmov, heuristík, metaheuristík aj moderných prístupov založených na strojovom učení.

Takto formulovaný problém možno interpretovať ako úlohu na úplnom grafe $G=(V,E)$, v ktorom vrcholy reprezentujú mestá a hrany náklady alebo vzdialenosti medzi nimi. Hľadaným riešením je hamiltonovský cyklus s minimálnou celkovou cenou. Táto grafová interpretácia je dôležitá nielen z hľadiska matematickej formulácie, ale aj z pohľadu návrhu algoritmov, pretože umožňuje priamo pracovať so štruktúrou vrcholov, hrán a susedností.

\begin{figure}[htbp]
    \centering
    \begin{tikzpicture}[
    city/.style={circle, fill=black, inner sep=2.2pt},
    route/.style={thick},
    alt/.style={gray!60, thin, dashed}
]

\coordinate (A) at (0.2,1.1);
\coordinate (B) at (1.3,2.7);
\coordinate (C) at (3.1,3.0);
\coordinate (D) at (4.8,2.2);
\coordinate (E) at (4.2,0.5);
\coordinate (F) at (2.4,0.2);
\coordinate (G) at (1.0,0.4);

% pomocné spojenia
\draw[alt] (A)--(C);
\draw[alt] (A)--(D);
\draw[alt] (B)--(D);
\draw[alt] (B)--(E);
\draw[alt] (C)--(E);
\draw[alt] (C)--(F);
\draw[alt] (D)--(F);
\draw[alt] (D)--(G);
\draw[alt] (E)--(G);
\draw[alt] (B)--(F);

% výsledná trasa
\draw[route] (A)--(B)--(C)--(D)--(E)--(F)--(G)--(A);

% vrcholy
\node[city,label=below left:$v_1$] at (A) {};
\node[city,label=above:$v_2$] at (B) {};
\node[city,label=above:$v_3$] at (C) {};
\node[city,label=above right:$v_4$] at (D) {};
\node[city,label=below right:$v_5$] at (E) {};
\node[city,label=below:$v_6$] at (F) {};
\node[city,label=below left:$v_7$] at (G) {};

\end{tikzpicture}
    \caption{Ilustračná euklidovská inštancia symetrického problému obchodného cestujúceho s vyznačenou výslednou trasou.}
    \label{fig:tsp_instance}
\end{figure}

Z tejto základnej formulácie vychádza viacero variantov problému TSP, ktoré sa odlišujú najmä vlastnosťami nákladovej matice, počtom obchodných cestujúcich alebo prítomnosťou dodatočných obmedzení. Medzi najčastejšie uvádzané varianty patria:

\begin{itemize}[leftmargin=1.5cm,itemsep=3pt,topsep=4pt]
    \item \textbf{STSP (Symmetric Traveling Salesman Problem),} pri ktorom je náklad presunu z vrcholu $i$ do vrcholu $j$ rovnaký ako náklad v opačnom smere;
    \item \textbf{ATSP (Asymmetric Traveling Salesman Problem),} v ktorom sa náklady presunu môžu líšiť podľa smeru pohybu;
    \item \textbf{mTSP (Multiple Traveling Salesman Problem),} kde je úloha rozdelená medzi viacerých obchodných cestujúcich;
    \item \textbf{TSPTW (Traveling Salesman Problem with Time Windows),} pri ktorom musia byť jednotlivé vrcholy navštívené v určených časových intervaloch;
    \item \textbf{viackriteriálne varianty TSP,} v ktorých sa optimalizácia nevzťahuje iba na jednu cieľovú funkciu, ale súčasne napríklad na čas, náklady alebo riziko.
\end{itemize}

Jednotlivé varianty vznikli ako reakcia na praktické požiadavky v logistike, plánovaní výroby, skladovaní, rozvoze alebo v navigačných systémoch. Hoci sa medzi sebou líšia konkrétnymi obmedzeniami, všetky vychádzajú z rovnakého princípu: nájsť efektívne poradie návštevy viacerých bodov pri rešpektovaní definovaných pravidiel úlohy.

V tejto práci sa zameriavame na symetrický problém obchodného cestujúceho. V tomto variante platí, že náklad presunu z vrcholu $i$ do vrcholu $j$ je rovnaký ako náklad presunu z vrcholu $j$ do vrcholu $i$. STSP predstavuje klasickú a metodologicky dobre preskúmanú formuláciu, ktorá sa často používa ako základ pri porovnávaní algoritmov, benchmarkových inštancií a metrík hodnotenia. Zároveň ide o variant, ktorý je vhodný aj pre experimentálne porovnávanie dátovo riadených modelov, keďže umožňuje pracovať s jednotnou a dobre interpretovateľnou štruktúrou úlohy.

Význam TSP však nespočíva iba v jeho teoretickej elegancii. Problém má široké praktické uplatnenie v logistike, distribúcii, plánovaní servisných trás, sekvenovaní výrobných operácií, návrhu dosiek plošných spojov, bioinformatike či v konštrukcii testovacích úloh pre nové optimalizačné postupy. Práve spojenie jasnej matematickej definície, praktickej relevancie a vysokej výpočtovej náročnosti vysvetľuje, prečo TSP zostáva dlhodobo jednou z centrálnych úloh operačného výskumu.

\section{Výpočtová náročnosť a klasické prístupy k riešeniu}

Na riešenie problému obchodného cestujúceho bolo navrhnuté široké spektrum prístupov, ktoré sa postupne vyvíjali v závislosti od požiadaviek na presnosť, rýchlosť a praktickú použiteľnosť. V tejto časti sa zameriavame na vysvetlenie výpočtovej náročnosti problému a na hlavné skupiny klasických metód používaných pri jeho riešení.

\noindent\textbf{Výpočtová náročnosť problému.}

Problém obchodného cestujúceho patrí medzi výpočtovo náročné optimalizačné úlohy. Jednoducho povedané to znamená, že počet možných trás rastie s počtom vrcholov veľmi rýchlo, a preto už pri stredne veľkých inštanciách nie je možné ani efektívne, ani časovo prijateľné úplne prehľadať priestor všetkých riešení. V prípade symetrického TSP je počet rôznych cyklov rádovo $(n-1)!/2$, čo dobre ilustruje, prečo sa problém s rastúcim počtom miest stáva prakticky neriešiteľným úplným prehľadaním priestoru riešení.

Práve táto vlastnosť mala zásadný vplyv na vývoj metód riešenia. Ukázalo sa, že pri menších inštanciách možno hľadať optimálne riešenie pomocou presných algoritmov, no pri väčších úlohách je potrebné pracovať s prístupmi, ktoré uprednostňujú rýchlosť a praktickú použiteľnosť pred garanciou absolútneho optima. Z historického aj metodologického hľadiska sa preto vyprofilovali dve hlavné skupiny klasických prístupov: presné metódy a heuristické, resp. metaheuristické postupy.

\noindent\textbf{Presné metódy riešenia.}

Presné metódy sú zamerané na nájdenie garantovane optimálneho riešenia. Medzi najdôležitejšie prístupy v tejto skupine patria celočíselné lineárne formulácie, vetvenie a medze, vetvenie s rezmi či dynamické programovanie. Ich spoločným znakom je to, že pri splnení výpočtových podmienok umožňujú určiť skutočne najlepšiu možnú trasu, a preto predstavujú dôležitý referenčný základ pri hodnotení kvality iných algoritmov.

Význam presných metód nespočíva iba v samotnom riešení menších alebo stredne veľkých inštancií. Ich výsledky zohrávajú dôležitú úlohu aj pri benchmarkingu, pretože poskytujú referenčné hodnoty, voči ktorým možno porovnávať približné a dátovo riadené prístupy. Zároveň slúžia ako zdroj kvalitných štítkov pri tvorbe trénovacích množín, keďže práve optimálne alebo takmer optimálne riešenia umožňujú učiť modely na spoľahlivých cieľových výstupoch.

Nevýhodou presných metód je ich vysoká výpočtová náročnosť, ktorá sa prejavuje najmä pri rastúcej veľkosti inštancie. V praxi sa preto využívajú najmä tam, kde je potrebná garancia optimálnosti, alebo tam, kde je možné akceptovať vyššie časové a výpočtové náklady výmenou za maximálnu presnosť výsledku.

\noindent\textbf{Heuristiky a metaheuristiky.}

V situáciách, keď presné riešenie nie je časovo alebo výpočtovo realistické, sa uplatňujú heuristiky a metaheuristiky. Heuristické metódy konštruujú riešenie pomocou relatívne jednoduchých rozhodovacích pravidiel, napríklad postupným výberom najbližšieho suseda alebo vkladaním ďalších vrcholov do už vytváranej trasy. Ich hlavnou výhodou je rýchlosť a implementačná jednoduchosť, no spravidla neposkytujú garanciu optimálnosti a bývajú citlivé na lokálne rozhodnutia urobené v počiatočných krokoch.

Metaheuristiky predstavujú všeobecnejší rámec, v ktorom sa kombinuje konštrukcia riešenia s mechanizmami zlepšovania, diverzifikácie a vyhľadávania v priestore alternatívnych trás. Patria sem napríklad simulované žíhanie, tabu search, genetické algoritmy, mravčie kolónie alebo iteratívne lokálne prehľadávanie. Ich cieľom nie je nájsť garantované optimum, ale dosiahnuť čo najkvalitnejšie riešenie v prijateľnom čase.

V kontexte symetrického TSP zohrávajú osobitnú úlohu lokálne zlepšovacie techniky typu $k$-opt a heuristika Lin--Kernighan. Práve tieto prístupy ukázali, že aj bez formálnej garancie optimálnosti možno v praxi dosahovať veľmi kvalitné riešenia, ktoré sú často dostatočné pre reálne aplikácie aj pre experimentálne porovnávanie algoritmov.

\noindent\textbf{Význam klasických prístupov pre súčasný výskum.}

Klasické metódy riešenia TSP dnes nepredstavujú iba historický základ výskumu, ale stále zohrávajú dôležitú úlohu aj v súčasných prístupoch. Presné algoritmy poskytujú referenčné riešenia a heuristiky a metaheuristiky ostávajú praktickými nástrojmi na riešenie väčších inštancií. Navyše sa ukazuje, že mnohé moderné dátovo riadené a neurónové modely tieto klasické postupy nenahrádzajú úplne, ale skôr na ne nadväzujú.

V hybridných riešeniach sa napríklad model strojového učenia používa na predikciu sľubných hrán alebo na usmernenie prehľadávania, zatiaľ čo samotná výsledná trasa vzniká až pomocou heuristického dekódovania alebo následného lokálneho dolaďovania. Z toho vyplýva, že pochopenie klasických prístupov je nevyhnutné nielen pre teoretické uchopenie problému, ale aj pre správnu interpretáciu moderných výsledkov. Práve na tento vývoj prirodzene nadväzuje nasledujúca časť práce, ktorá sa venuje dátovo riadeným a neurónovým prístupom k riešeniu TSP.

\begin{figure}[htbp]
    \centering
    \begin{tikzpicture}[
    node distance=1.4cm and 1.6cm,
    box/.style={
        draw,
        rounded corners,
        align=center,
        minimum width=3.2cm,
        minimum height=0.95cm
    },
    arrow/.style={-{Latex[length=2.5mm]}, thick}
]

\node[box] (tsp) {Prístupy k riešeniu TSP};

\node[box, below left=of tsp] (exact) {Presné\\ metódy};
\node[box, below=of tsp] (heur) {Heuristiky a\\ metaheuristiky};
\node[box, below right=of tsp] (ml) {Dátovo riadené\\ prístupy};
\node[box, below=of heur] (hybrid) {Hybridné\\ prístupy};

\draw[arrow] (tsp) -- (exact);
\draw[arrow] (tsp) -- (heur);
\draw[arrow] (tsp) -- (ml);
\draw[arrow] (heur) -- (hybrid);
\draw[arrow] (ml) -- (hybrid);

\end{tikzpicture}
    \caption{Prehľad hlavných skupín prístupov k riešeniu problému obchodného cestujúceho.}
    \label{fig:tsp_methods}
\end{figure}

\section{Formulácia problému TSP pre dátovo riadené modely}

Pri klasickom riešení problému obchodného cestujúceho sa optimálna trasa hľadá priamo pomocou optimalizačného algoritmu. V prípade dátovo riadených prístupov je však potrebné túto úlohu najskôr previesť do podoby, s ktorou môže pracovať model strojového učenia. To znamená, že problém už nie je chápaný iba ako hľadanie najlepšej trasy, ale aj ako úloha nad dátami, v ktorej treba definovať vstupné inštancie, cieľové výstupy a vhodnú reprezentáciu riešenia.

Vstup do modelu zvyčajne tvorí množina vrcholov reprezentovaných súradnicami, maticou vzdialeností alebo odvodenými príznakmi hrán. Samotné vstupné inštancie však na učenie nestačia. Model potrebuje aj referenčné riešenia, na základe ktorých sa vytvoria cieľové štítky. V prípade TSP môže mať takýto štítok podobu celej optimálnej alebo približne optimálnej trasy, prípadne informácie o tom, ktoré hrany patria do referenčného riešenia.

Pri dátovo riadených prístupoch sa preto často využívajú synteticky generované inštancie, najmä euklidovské úlohy, pri ktorých sú vrcholy vytvorené náhodne v dvojrozmernom priestore. Výhodou takéhoto postupu je možnosť pripraviť veľké množstvo tréningových príkladov s kontrolovanou veľkosťou a štruktúrou. Zároveň však vzniká otázka kvality referenčných riešení. Presné riešiče poskytujú kvalitné štítky, ich použitie je však časovo a výpočtovo náročné. Aproximačné postupy sú rýchlejšie, no model sa pri nich učí z riešení, ktoré nemusia byť optimálne. Voľba medzi presnosťou štítkov a náročnosťou ich prípravy preto predstavuje dôležitý metodologický aspekt.

Pri návrhu modelu možno TSP formulovať viacerými spôsobmi. Jednou z možností je sekvenčné generovanie poradia vrcholov. Druhou, z pohľadu našej práce kľúčovou možnosťou, je formulácia problému ako predikcie hrán v úplnom grafe. V tomto prípade model nepredikuje priamo celé poradie miest, ale odhaduje, ktoré hrany sú dôležité pre výslednú trasu. Takáto formulácia je prirodzená preto, že samotný TSP je definovaný na grafe a riešenie možno chápať ako množinu hrán tvoriacich hamiltonovský cyklus.

Výhodou predikcie hrán je, že umožňuje pracovať s grafovou štruktúrou úlohy a s príznakmi odvodenými zo vzdialeností, geometrie alebo lokálneho usporiadania vrcholov. Z pohľadu učenia sa každá hrana stáva samostatnou jednotkou predikcie, pre ktorú model odhaduje skóre, pravdepodobnosť alebo triedu. Výstupom ešte nie je finálna trasa, ale štruktúrovaná informácia o tom, ktoré spojenia sú pre kvalitné riešenie perspektívne.

Dôležitým špecifikom tohto prístupu je, že samotná predikcia ešte nezaručuje vznik realizovateľného riešenia. Model môže priradiť vysoké skóre hranám, ktoré spolu nevytvárajú uzavretý cyklus spĺňajúci všetky podmienky TSP. Preto je medzi výstupom modelu a finálnou trasou nevyhnutná dekódovacia fáza, ktorá prevedie predikované skóre na konkrétny hamiltonovský cyklus. Táto fáza môže byť doplnená aj o lokálne zlepšovanie, čo vedie k hybridnému pipeline spájajúcemu dátovo riadený model s klasickými algoritmickými postupmi.

Pre potreby tejto práce má takáto formulácia zásadný význam, keďže umožňuje prepájať grafovú štruktúru TSP, prípravu tréningových dát a výstup neurónových modelov v jednotnom experimentálnom rámci. Práve na tomto základe následne stojí aj nasledujúca časť kapitoly, ktorá sa venuje konkrétnym dátovo riadeným a neurónovým prístupom k riešeniu TSP.

\section{Dátovo riadené a neurónové prístupy k riešeniu TSP}

V súčasnom výskume kombinatorickej optimalizácie sa čoraz väčšia pozornosť venuje prístupom, ktoré využívajú strojové učenie na odhad rozhodnutí vedúcich ku kvalitným riešeniam. Pri probléme obchodného cestujúceho to znamená, že model sa neučí priamo „riešiť“ úlohu v tradičnom zmysle, ale odhaduje štruktúrne vzory, ktoré možno následne využiť pri konštrukcii výslednej trasy. V našej práci sa zameriavame na taký typ dátovo riadených modelov, ktorý pracuje s hranami úplného grafu a hodnotí ich význam pre výsledné riešenie.

\noindent\textbf{Edge-based prístup ako základ modelovania.}

Pri edge-based formulácii model nepredikuje celé poradie miest, ale priraďuje skóre jednotlivým hranám medzi vrcholmi. Takýto prístup prirodzene vychádza z grafovej povahy problému TSP, keďže výslednú trasu možno chápať ako špecifickú množinu hrán tvoriacich hamiltonovský cyklus. Výhodou tejto formulácie je, že umožňuje pracovať s explicitne pripravenými príznakmi hrán a zároveň porovnávať rôzne architektúry v jednotnom experimentálnom rámci.

Z metodologického hľadiska ide o úlohu, v ktorej model odhaduje, ktoré spojenia medzi vrcholmi sú perspektívne pre výslednú trasu. Výstupom preto nie je priamo hotový okruh, ale skóre alebo pravdepodobnosť priradená jednotlivým hranám. Takáto reprezentácia je vhodná najmä vtedy, keď chceme analyzovať, do akej miery samotná architektúra modelu ovplyvňuje kvalitu predikcie nad rovnakou množinou vstupných príznakov.

\noindent\textbf{Viacvrstvové perceptróny nad príznakmi hrán.}

Najjednoduchšiu skupinu modelov v tomto rámci predstavujú viacvrstvové perceptróny, ktoré pracujú s pripravenými príznakmi hrán. Tieto modely posudzujú každú hranu na základe jej vstupnej reprezentácie a snažia sa odhadnúť, či je dané spojenie dôležité pre kvalitnú trasu. Ich výhodou je architektonická jednoduchosť, relatívne nízka výpočtová náročnosť a dobrá interpretovateľnosť. Zároveň poskytujú prirodzený základ na porovnanie s výkonnejšími architektúrami.

V praktickej rovine sú takéto modely vhodné najmä ako baseline, pretože umožňujú relatívne priamo sledovať vplyv počtu vrstiev a šírky siete na výslednú kvalitu predikcie. Ak model pracuje iba s pripravenými príznakmi jednotlivých hrán, jeho schopnosť zachytiť zložitejšie globálne vzťahy v grafe je obmedzená, no práve preto predstavuje dôležitý referenčný bod pri porovnávaní s hlbšími a štruktúrne bohatšími architektúrami.

\noindent\textbf{Reziduálne siete pre hlbšiu reprezentáciu hrán.}

Ďalšiu skupinu predstavujú reziduálne siete, ktoré rozširujú jednoduchšie MLP modely o reziduálne prepojenia medzi vrstvami. Ich význam spočíva najmä v tom, že umožňujú stabilnejšie učenie hlbších architektúr a znižujú riziko degradácie výkonu pri väčšom počte vrstiev. V kontexte predikcie hrán to znamená, že model môže vytvárať bohatšiu vnútornú reprezentáciu bez toho, aby sa tréning výrazne destabilizoval.

Reziduálny prístup je preto vhodný najmä v situáciách, keď jednoduchý perceptrón už nedokáže dostatočne zachytiť zložitejšie vzťahy medzi vstupnými príznakmi. Takýto model síce stále pracuje primárne nad edge features, no vďaka hlbšej transformácii môže lepšie využívať ich kombinácie a nelineárne vzťahy. V našom prípade preto predstavuje prirodzený medzistupeň medzi jednoduchým baseline modelom a štruktúrne uvedomelejším transformerovým prístupom.

\noindent\textbf{Transformerové modely so zohľadnením štruktúry grafu.}

Najbohatšiu architektonickú skupinu v našej práci predstavujú modely využívajúce mechanizmy pozornosti. Na rozdiel od jednoduchších modelov, ktoré pracujú iba s izolovanými príznakmi hrán, tieto architektúry dokážu zohľadňovať aj širšie vzťahy medzi vrcholmi a hranami v celej inštancii. V našom prípade ide o prístup, v ktorom sa najskôr vytvára reprezentácia vrcholov na základe ich súradníc a následne sa skórujú jednotlivé hrany s využitím edge features aj kontextu odvodeného z uzlových reprezentácií.

Takýto prístup je dôležitý najmä preto, že kvalita konkrétnej hrany v TSP nezávisí iba od jej lokálnych geometrických vlastností, ale aj od globálneho usporiadania celej inštancie. Transformerový model preto nepracuje len s lokálnym popisom hrany, ale zároveň zohľadňuje širšiu štruktúru grafu. Práve tým sa odlišuje od jednoduchších MLP a reziduálnych architektúr a predstavuje prirodzený krok smerom k grafovo uvedomelému modelovaniu.

\begin{figure}[htbp]
    \centering
    \begin{tikzpicture}[
    node distance=0.85cm and 0.65cm,
    box/.style={
        draw,
        rounded corners,
        align=center,
        minimum width=1.95cm,
        minimum height=0.82cm,
        inner sep=3pt
    },
    arrow/.style={-{Latex[length=2mm]}, thick}
]

% horný riadok
\node[box] (input) {Vstupná\\ inštancia};
\node[box, right=of input] (graph) {Úplný\\ graf};
\node[box, right=of graph] (features) {Príznaky\\ hrán};
\node[box, right=of features] (model) {Neurónový\\ model};
\node[box, right=of model] (scores) {Skóre\\ hrán};

% dolný riadok – kompaktnejší a centrovaný
\node[box, below=1.1cm of model] (decode) {Dekódovanie};
\node[box, left=of decode] (opt) {2-opt / lokálne\\ zlepšenie};
\node[box, left=of opt] (route) {Finálna\\ trasa};

% šípky horný riadok
\draw[arrow] (input) -- (graph);
\draw[arrow] (graph) -- (features);
\draw[arrow] (features) -- (model);
\draw[arrow] (model) -- (scores);

% prechod dole a dokončenie pipeline
\draw[arrow] (scores) |- (decode);
\draw[arrow] (decode) -- (opt);
\draw[arrow] (opt) -- (route);

\end{tikzpicture}
    \caption{Schéma edge-based pipeline od vstupnej inštancie cez výpočet príznakov hrán a predikciu skóre až po dekódovanie výslednej trasy.}
    \label{fig:edge_pipeline}
\end{figure}

\noindent\textbf{Dekódovanie a lokálna optimalizácia.}

Spoločnou črtou uvedených modelov je, že ich výstupom nie je finálna realizovateľná trasa, ale iba skóre hrán. Z tohto dôvodu je nevyhnutná dekódovacia fáza, ktorá prevedie predikované hodnoty na konkrétny hamiltonovský cyklus. Bez tejto fázy by model mohol označiť ako perspektívne také hrany, ktoré spolu nevytvárajú platné riešenie.

V praxi sa preto dátovo riadený model prepája s klasickými algoritmickými postupmi. Dekódovanie býva často doplnené aj o lokálne zlepšovanie, napríklad pomocou techník typu 2-opt, ktoré odstraňujú zjavne neefektívne prepojenia a zlepšujú výslednú trasu. Moderné neurónové prístupy k TSP tak zvyčajne nefungujú ako úplná náhrada klasických riešičov, ale ako súčasť hybridného pipeline, v ktorom sa spája predikcia založená na dátach s algoritmickou konštrukciou a dolaďovaním riešenia.

\section{Benchmarkové dáta a metriky hodnotenia}

\noindent\textbf{Referenčné benchmarky a syntetické dáta.}

Objektívne porovnávanie metód riešenia TSP si vyžaduje štandardizované benchmarky. Za najvýznamnejší referenčný zdroj sa v tejto oblasti považuje knižnica TSPLIB, ktorá obsahuje široké spektrum symetrických aj asymetrických inštancií rôznych veľkostí a rôznych metrík vzdialenosti. TSPLIB predstavuje dlhodobo používaný etalón pri hodnotení kvality algoritmov a umožňuje porovnávať nové prístupy s výsledkami uvádzanými v odbornej literatúre.

Jej význam nespočíva iba v samotnej dostupnosti testovacích inštancií, ale aj v tom, že pri mnohých úlohách sú k dispozícii optimálne alebo najlepšie známe referenčné riešenia. Práve táto vlastnosť robí z TSPLIB mimoriadne dôležitý nástroj na hodnotenie kvality navrhnutých modelov, pretože umožňuje porovnávať výsledné trasy voči spoľahlivému referenčnému štandardu. V kontexte tejto práce je to obzvlášť dôležité, keďže kvalita modelov nie je posudzovaná iba podľa samotnej predikcie hrán, ale najmä podľa kvality výslednej dekódovanej trasy.

Pre dátovo riadené prístupy však samotná TSPLIB zvyčajne nestačí ako tréningový zdroj, pretože obsahuje relatívne obmedzený počet inštancií. Z tohto dôvodu sa pri učení modelov často využívajú synteticky generované euklidovské dáta, na ktorých možno pripraviť veľké množstvo označených príkladov s kontrolovanou distribúciou vrcholov a veľkosťou grafu. Zatiaľ čo syntetické inštancie slúžia predovšetkým na tréning a experimentovanie, TSPLIB plní úlohu nezávislého benchmarku na overenie generalizačnej schopnosti modelu a na porovnanie s existujúcimi riešeniami.

Pri benchmarkoch je zároveň nevyhnutné rešpektovať konkrétnu metriku danej inštancie. TSPLIB obsahuje napríklad úlohy s metrikami \texttt{EUC\_2D}, \texttt{CEIL\_2D}, \texttt{ATT} alebo \texttt{GEO}, pričom každá z nich má vlastné pravidlá výpočtu dĺžky hrany. Pre korektné porovnanie výsledkov preto nestačí pracovať iba s normalizovanými súradnicami; konečné hodnotenie musí vychádzať z metricky správne vypočítanej dĺžky trasy podľa definície konkrétnej inštancie.

\begin{figure}[htbp]
    \centering
    \begin{tikzpicture}[
    node distance=1.5cm and 1.3cm,
    box/.style={
        draw,
        rounded corners,
        align=center,
        minimum width=3.3cm,
        minimum height=1cm
    },
    arrow/.style={-{Latex[length=2.5mm]}, thick}
]

\node[box] (data) {Dátové zdroje\\ pre hodnotenie};

\node[box, below left=of data] (synthetic) {Syntetické\\ inštancie};
\node[box, below right=of data] (tsplib) {Inštancie\\ TSPLIB};

\node[box, below=of synthetic] (train) {Tréning a\\ experimentovanie};
\node[box, below=of tsplib] (eval) {Nezávislé\\ hodnotenie};

\node[box, below=1.7cm of $(train)!0.5!(eval)$] (ref) {Referenčné riešenia\\ a relatívna odchýlka};

\draw[arrow] (data) -- (synthetic);
\draw[arrow] (data) -- (tsplib);
\draw[arrow] (synthetic) -- (train);
\draw[arrow] (tsplib) -- (eval);
\draw[arrow] (train) |- (ref);
\draw[arrow] (eval) |- (ref);

\end{tikzpicture}

    \caption{Prehľad hlavných zdrojov dát používaných pri tréningu a hodnotení modelov pre TSP.}
    \label{fig:benchmarks}
\end{figure}

\noindent\textbf{Metriky porovnávania riešení.}

Najzákladnejšou metrikou kvality riešenia je dĺžka výslednej trasy. Pri porovnávaní neurónových a heuristických prístupov sa však často používa aj optimalizačná odchýlka (optimality gap), ktorá vyjadruje relatívny rozdiel medzi dĺžkou predikovanej trasy a referenčnou hodnotou. Túto metriku možno zapísať v tvare
\begin{equation}
\mathrm{gap}(\%) = \frac{L_{\mathrm{pred}} - L_{\mathrm{ref}}}{L_{\mathrm{ref}}} \cdot 100,
\end{equation}
pričom $L_{\mathrm{ref}}$ zodpovedá optimálnej alebo najlepšej známej referenčnej trase.

Výhodou optimalizačnej odchýlky je, že umožňuje zmysluplne porovnávať inštancie rôznych veľkostí a rôznych absolútnych dĺžok trás. Okrem priemernej odchýlky sa pri hodnotení sledujú aj ďalšie charakteristiky, napríklad medián, percentily, stabilita medzi opakovaniami, podiel platných riešení alebo čas inferencie. Pri edge-based prístupoch je navyše dôležité posudzovať nielen samotné skóre hrán, ale najmä kvalitu výslednej trasy po dekódovaní a prípadnej lokálnej optimalizácii.

\section{Zhodnotenie súčasného stavu a výskumná medzera}

Z prehľadu súčasného stavu riešenej problematiky vyplýva, že výskum TSP sa dnes opiera o tri vzájomne prepojené línie. Prvú tvoria presné metódy a silné heuristiky, ktoré zostávajú referenčným základom pri hodnotení kvality riešení. Druhú predstavuje formulácia problému pre dátovo riadené modely, v ktorej zohrávajú dôležitú úlohu syntetické inštancie, referenčné riešenia a vhodne definované cieľové štítky. Tretiu líniu tvoria samotné neurónové architektúry, ktoré sa snažia využiť pripravené dáta na odhad štrukturálnych vzorov vedúcich ku kvalitným trasám.

Napriek výraznému rozvoju tejto oblasti zostáva viacero otvorených otázok. Naďalej nie je jednoznačné, aký typ edge-based architektúry poskytuje najlepší kompromis medzi jednoduchosťou modelu, schopnosťou zachytiť štruktúru grafu a kvalitou výsledného riešenia. Rovnako dôležitou otázkou zostáva vplyv kvality referenčných štítkov a použitých dekódovacích postupov na finálny výkon modelu. Významnou výzvou je aj to, do akej miery sa výsledky dosiahnuté na syntetických dátach prenášajú na štandardné benchmarky typu TSPLIB.

Za výskumnú medzeru preto možno považovať potrebu jednotného experimentálneho rámca, v ktorom by bolo možné transparentne porovnať viacero edge-based modelov pre symetrický TSP na rovnakých dátach, s rovnakou reprezentáciou príznakov, s porovnateľným dekódovaním a s jednotnou metrikou hodnotenia. Práve na túto medzeru nadväzuje aj naša práca. V nasledujúcej kapitole preto prejdeme od všeobecného teoretického rámca ku konkrétnej metodike, v ktorej opíšeme spôsob prípravy dát, návrh modelov, dekódovanie riešení a experimentálne nastavenie hodnotenia.
\chapter{Cieľ práce, metodika práce a metódy skúmania}

%% ═══════════════════════════════════════════════════════════════════════════
%% 3.1  CIEĽ PRÁCE
%% ═══════════════════════════════════════════════════════════════════════════
\section{Cieľ práce}
\label{sec:ciel}

\indent Z prehľadu súčasného stavu vyplynulo, že neurónovým prístupom
k problému obchodného cestujúceho chýba spoločný metodický základ:
práce sa líšia formuláciou úlohy, kvalitou referenčných riešení aj
spôsobom hodnotenia, čo znemožňuje priame porovnanie architektúr.
Práve táto medzera motivovala zvolený prístup — grafovú formuláciu
postavenou na predikcii hrán s jednotným hodnotením na štandardných
inštanciách TSPLIB
\parencite{joshi2022generalization,bengio2021mlco,reinelt1991tsplib}.

\AssignmentGoalSK

\indent Tretiu zo stanovených úloh sme v rámci práce rozvinuli do
piatich merateľných čiastkových cieľov, ktoré pokrývajú celý postup
od prípravy údajov až po záverečnú interpretáciu výsledkov:

\begin{enumerate}
    \item Pripraviť syntetické tréningové a validačné údaje s dvoma
          úrovňami kvality referenčných riešení a spracovať štandardné
          testovacie inštancie z knižnice TSPLIB.
    \item Navrhnúť a implementovať tri rodiny modelov založených na
          predikcii hrán v úplnom grafe, od jednoduchého hranového
          perceptrónu cez reziduálny variant až po kontextový transformer.
    \item Vytvoriť jednotný hodnotiaci rámec, v ktorom sa modely aj
          heuristické bázy porovnávajú na rovnakých štrnástich inštanciách
          TSPLIB za rovnakých dekódovacích podmienok.
    \item Porovnať vplyv kvality tréningových štítkov a dekódovacích
          stratégií na výslednú kvalitu riešení.
    \item Analyzovať, či kontextový model prináša merateľné zlepšenie
          oproti lokálnym hranovým modelom, a identifikovať podmienky,
          v ktorých sú jednotlivé prístupy najefektívnejšie.
\end{enumerate}

%% ═══════════════════════════════════════════════════════════════════════════
%% 3.2  PREHĽAD EXPERIMENTÁLNEHO POSTUPU
%% ═══════════════════════════════════════════════════════════════════════════
\section{Prehľad experimentálneho postupu}
\label{sec:prehlad_postupu}

\indent Praktická časť práce bola realizovaná ako ucelený experimentálny
postup, v ktorom jednotlivé kroky priamo nadväzujú jeden na druhý. Keďže
pri dátovo riadených prístupoch ku kombinatorickej optimalizácii nemožno
oddeliť model od spôsobu prípravy dát, kvality referenčných riešení
a následného spracovania výstupu, metodická časť je koncipovaná ako opis
tohto celku, nie len jeho modelovej zložky
\parencite{bengio2021mlco,joshi2022generalization,cappart2023gnnco}.

\indent Z hľadiska údajov bol postup rozdelený na dve vzájomne prepojené
vetvy. Prvú tvorili synteticky generované euklidovské inštancie, nad
ktorými prebiehalo trénovanie a validácia modelov, pričom referenčné
riešenia sa pripravili v dvoch variantoch kvality: heuristickom
a presnom pomocou riešiča Concorde \parencite{concorde2003solver}. Druhú
vetvu tvorili štandardné inštancie z knižnice TSPLIB, ktoré slúžili
výhradne na nezávislé záverečné hodnotenie a neboli súčasťou trénovacej
množiny \parencite{reinelt1991tsplib}. Toto oddelenie umožnilo sledovať,
do akej miery sa naučené vzory prenášajú na úlohy odlišné od synteticky
generovaných príkladov.

\indent Nad oboma vetvami sa uplatnila jednotná modelová časť: všetky
inštancie sa previedli na grafovú reprezentáciu a modely sa učili
predikovať príslušnosť hrán k výslednej trase. Takáto formulácia
udržiava rovnaké experimentálne podmienky pre všetky tri porovnávané
rodiny architektúr a priamo previazuje modelový výstup s následným
dekódovaním trasy. Celkový návrh zachytáva obrázok~\ref{fig:method_overview}.

\begin{figure}[H]
    \centering
    \resizebox{0.96\textwidth}{!}{% method_overview.tex
% Vyžaduje \usetikzlibrary{arrows.meta,positioning,calc}
% Použitie: \resizebox{0.90\textwidth}{!}{% method_overview.tex
% Vyžaduje \usetikzlibrary{arrows.meta,positioning,calc}
% Použitie: \resizebox{0.90\textwidth}{!}{\input{figures/tikz/method_overview}}
\begin{tikzpicture}[
    font=\small,
    node distance=0.65cm,
    header/.style={
        draw, rounded corners=4pt, align=center,
        minimum width=3.7cm, inner sep=5pt,
        font=\small\bfseries,
        text=#1!80!black, fill=#1!10, draw=#1!55
    },
    train/.style={
        draw=blue!40, fill=blue!6, rounded corners=3pt,
        align=center, minimum width=3.7cm, inner sep=5pt
    },
    refbox/.style={
        draw=orange!45, fill=orange!8, rounded corners=3pt,
        align=center, minimum width=1.72cm, inner sep=5pt
    },
    evalbox/.style={
        draw=green!50, fill=green!6, rounded corners=3pt,
        align=center, minimum width=3.7cm, inner sep=5pt
    },
    decodebox/.style={
        draw=gray!50, fill=gray!8, rounded corners=3pt,
        align=center, minimum width=5.4cm, inner sep=5pt
    },
    finalbox/.style={
        draw=gray!55, fill=gray!12, rounded corners=4pt,
        align=center, minimum width=9.8cm, inner sep=5pt
    },
    arrow/.style={
        -{Latex[length=2.6mm, width=1.9mm]}, thick,
        shorten >=3pt
    },
    darrow/.style={
        -{Latex[length=2.6mm, width=1.9mm]}, thick,
        dashed, draw=gray!60, shorten >=3pt
    },
    pline/.style={thick}
]

%% ── Uzly ──────────────────────────────────────────────────────────
\node[header=blue]  (htrain) at (0.0, 0.0) {Tréningová vetva};
\node[header=green] (heval)  at (6.8, 0.0) {Hodnotiaca vetva};

\node[train, below=of htrain] (synth)
    {Syntetické inštancie\\[2pt]\(n \in \{20,30,\ldots,100\}\)};

\node[train, below=of synth] (ref)
    {Referenčné riešenia};

%% heur a conc: väčší gap 1.40cm — dostatok priestoru pre šípky
\node[refbox, below=1.40cm of ref, xshift=-1.18cm] (heur)
    {Heuristické\\[2pt]\scriptsize(NN + 2-opt)};
\node[refbox, below=1.40cm of ref, xshift= 1.18cm] (conc)
    {Presné\\[2pt]\scriptsize(Concorde)};

%% trainmod: štandardný gap od spodku heur/conc
\node[train, below=1.20cm of heur.south |- conc.south] (trainmod)
    {Trénovanie modelov\\[2pt]MLP,\ ResMLP,\ Transformer};

\node[evalbox] at (heval |- synth)    (tsplib)
    {14 inštancií TSPLIB\\[2pt]štandardné benchmarky};
\node[evalbox] at (heval |- trainmod) (evalmod)
    {Hodnotenie modelov\\[2pt]na inštanciách TSPLIB};

\node[decodebox, below=0.90cm of trainmod] (decode)
    {Dekódovanie + lokálne zlepšenie (2-opt)};
\node[finalbox,  below=0.65cm of decode]   (anal)
    {Záverečná analýza a porovnanie výsledkov};

%% ── Šípky ─────────────────────────────────────────────────────────
\draw[arrow] (htrain.south) -- (synth.north);
\draw[arrow] (synth.south)  -- (ref.north);

%% Fork: fsplit = stred medzi ref.south a heur.north
%% pri gap=1.40cm: priestor = 1.40cm → fsplit 0.70cm od každého
\coordinate (fsplit) at ($(ref.south)!0.5!(ref.south |- heur.north)$);

\draw[pline] (ref.south) -- (fsplit);
\draw[arrow] (fsplit) -- (fsplit -| heur) -- (heur.north);
\draw[arrow] (fsplit) -- (fsplit -| conc) -- (conc.north);

%% Merge: fmerge = stred medzi heur.south a trainmod.north
\coordinate (fmerge) at ($(heur.south |- trainmod.north)!0.5!(heur.south)$);

\draw[pline] (heur.south) -- (heur.south |- fmerge) -- (fmerge);
\draw[pline] (conc.south) -- (conc.south |- fmerge) -- (fmerge);
\draw[arrow] (fmerge) -- (trainmod.north);

%% Pravá vetva
\draw[arrow] (heval.south)  -- (tsplib.north);
\draw[arrow] (tsplib.south) -- (evalmod.north);

%% Zdieľaná šípka
\draw[darrow] (trainmod.east) --
    node[above, font=\scriptsize\itshape, yshift=1pt]{zdieľané váhy}
    (evalmod.west);

%% Konvergencia → decode → anal
\coordinate (jbar) at ($(trainmod.south)!0.5!(trainmod.south |- decode.north)$);
\draw[pline] (trainmod.south) -- (trainmod.south |- jbar);
\draw[pline] (evalmod.south)  -- (evalmod.south  |- jbar);
\draw[pline] (trainmod.south |- jbar) -- (evalmod.south |- jbar);
\draw[arrow] (trainmod |- jbar) -- (decode.north);
\draw[arrow] (decode.south) -- (anal.north);

\end{tikzpicture}}
\begin{tikzpicture}[
    font=\small,
    node distance=0.65cm,
    header/.style={
        draw, rounded corners=4pt, align=center,
        minimum width=3.7cm, inner sep=5pt,
        font=\small\bfseries,
        text=#1!80!black, fill=#1!10, draw=#1!55
    },
    train/.style={
        draw=blue!40, fill=blue!6, rounded corners=3pt,
        align=center, minimum width=3.7cm, inner sep=5pt
    },
    refbox/.style={
        draw=orange!45, fill=orange!8, rounded corners=3pt,
        align=center, minimum width=1.72cm, inner sep=5pt
    },
    evalbox/.style={
        draw=green!50, fill=green!6, rounded corners=3pt,
        align=center, minimum width=3.7cm, inner sep=5pt
    },
    decodebox/.style={
        draw=gray!50, fill=gray!8, rounded corners=3pt,
        align=center, minimum width=5.4cm, inner sep=5pt
    },
    finalbox/.style={
        draw=gray!55, fill=gray!12, rounded corners=4pt,
        align=center, minimum width=9.8cm, inner sep=5pt
    },
    arrow/.style={
        -{Latex[length=2.6mm, width=1.9mm]}, thick,
        shorten >=3pt
    },
    darrow/.style={
        -{Latex[length=2.6mm, width=1.9mm]}, thick,
        dashed, draw=gray!60, shorten >=3pt
    },
    pline/.style={thick}
]

%% ── Uzly ──────────────────────────────────────────────────────────
\node[header=blue]  (htrain) at (0.0, 0.0) {Tréningová vetva};
\node[header=green] (heval)  at (6.8, 0.0) {Hodnotiaca vetva};

\node[train, below=of htrain] (synth)
    {Syntetické inštancie\\[2pt]\(n \in \{20,30,\ldots,100\}\)};

\node[train, below=of synth] (ref)
    {Referenčné riešenia};

%% heur a conc: väčší gap 1.40cm — dostatok priestoru pre šípky
\node[refbox, below=1.40cm of ref, xshift=-1.18cm] (heur)
    {Heuristické\\[2pt]\scriptsize(NN + 2-opt)};
\node[refbox, below=1.40cm of ref, xshift= 1.18cm] (conc)
    {Presné\\[2pt]\scriptsize(Concorde)};

%% trainmod: štandardný gap od spodku heur/conc
\node[train, below=1.20cm of heur.south |- conc.south] (trainmod)
    {Trénovanie modelov\\[2pt]MLP,\ ResMLP,\ Transformer};

\node[evalbox] at (heval |- synth)    (tsplib)
    {14 inštancií TSPLIB\\[2pt]štandardné benchmarky};
\node[evalbox] at (heval |- trainmod) (evalmod)
    {Hodnotenie modelov\\[2pt]na inštanciách TSPLIB};

\node[decodebox, below=0.90cm of trainmod] (decode)
    {Dekódovanie + lokálne zlepšenie (2-opt)};
\node[finalbox,  below=0.65cm of decode]   (anal)
    {Záverečná analýza a porovnanie výsledkov};

%% ── Šípky ─────────────────────────────────────────────────────────
\draw[arrow] (htrain.south) -- (synth.north);
\draw[arrow] (synth.south)  -- (ref.north);

%% Fork: fsplit = stred medzi ref.south a heur.north
%% pri gap=1.40cm: priestor = 1.40cm → fsplit 0.70cm od každého
\coordinate (fsplit) at ($(ref.south)!0.5!(ref.south |- heur.north)$);

\draw[pline] (ref.south) -- (fsplit);
\draw[arrow] (fsplit) -- (fsplit -| heur) -- (heur.north);
\draw[arrow] (fsplit) -- (fsplit -| conc) -- (conc.north);

%% Merge: fmerge = stred medzi heur.south a trainmod.north
\coordinate (fmerge) at ($(heur.south |- trainmod.north)!0.5!(heur.south)$);

\draw[pline] (heur.south) -- (heur.south |- fmerge) -- (fmerge);
\draw[pline] (conc.south) -- (conc.south |- fmerge) -- (fmerge);
\draw[arrow] (fmerge) -- (trainmod.north);

%% Pravá vetva
\draw[arrow] (heval.south)  -- (tsplib.north);
\draw[arrow] (tsplib.south) -- (evalmod.north);

%% Zdieľaná šípka
\draw[darrow] (trainmod.east) --
    node[above, font=\scriptsize\itshape, yshift=1pt]{zdieľané váhy}
    (evalmod.west);

%% Konvergencia → decode → anal
\coordinate (jbar) at ($(trainmod.south)!0.5!(trainmod.south |- decode.north)$);
\draw[pline] (trainmod.south) -- (trainmod.south |- jbar);
\draw[pline] (evalmod.south)  -- (evalmod.south  |- jbar);
\draw[pline] (trainmod.south |- jbar) -- (evalmod.south |- jbar);
\draw[arrow] (trainmod |- jbar) -- (decode.north);
\draw[arrow] (decode.south) -- (anal.north);

\end{tikzpicture}}
    \caption{Celkový návrh experimentálneho postupu: tréningová vetva
             (syntetické dáta a referenčné riešenia), vetva hodnotenia
             (inštancie TSPLIB) a záverečná analýza.}
    \label{fig:method_overview}
\end{figure}

\indent Implementácia bola postavená na jazyku Python s knižnicami NumPy,
PyTorch a Matplotlib. Konfigurácia každého experimentu sa riadila
súbormi YAML a výstupy sa ukladali vo formátoch \texttt{npz},
\texttt{JSON} a \texttt{CSV}. Každý krok má vlastný konfiguračný vstup
a vlastný súbor výsledkov, čo zaručuje opakovateľnosť a vylučuje
skryté manuálne zásahy medzi fázami.

%% ═══════════════════════════════════════════════════════════════════════════
%% 3.3  PRÍPRAVA ÚDAJOV
%% ═══════════════════════════════════════════════════════════════════════════
\section{Príprava údajov}
\label{sec:priprava_udajov}

\indent Táto sekcia podrobne opisuje prípravu oboch údajových vetiev
naznačených v predchádzajúcom prehľade: syntetickú tréningovú množinu
vrátane dvoch úrovní kvality referenčných riešení a štandardné testovacie
inštancie z knižnice TSPLIB.

%% ─── 3.3.1 Syntetické údaje ─────────────────────────────────────────────
\subsection{Syntetické tréningové a validačné údaje}
\label{sec:synteticke_udaje}

\indent Syntetická tréningová množina pozostávala z náhodne generovaných
euklidovských inštancií. Pre každú veľkosť \(n \in \{20, 30, 40, 50, 60,
70, 80, 90, 100\}\) sa vygenerovali tréningová, validačná a testovacia
podmnožina v pomere \(45\,000 : 9\,000 : 9\,000\) inštancií. Súradnice
miest každej inštancie sa vzorkovali nezávisle z rovnomerného rozdelenia
\(\mathcal{U}([0,1]^2)\).

\indent Keďže syntetická množina vyžadovala generovanie vo veľkom objeme,
bolo potrebné vyvážiť kvalitu referenčných riešení a výpočtovú náročnosť
ich prípravy. Z tohto dôvodu sa použili dve úrovne kvality.

\indent \textit{Približná heuristická úroveň.}
Základný konštrukčný postup vychádzal z heuristiky najbližšieho suseda.
Po výbere pevne určeného štartovacieho vrcholu sa počiatočná trasa
budovala postupne podľa vzťahu
\begin{equation}
\pi^{(0)}_1 = 1, \qquad
\pi^{(0)}_{t+1}
=
\arg\min_{j \in U_t} d_{\pi^{(0)}_t,j},
\qquad
U_t = V \setminus \{\pi^{(0)}_1,\dots,\pi^{(0)}_t\},
\label{eq:nearest_neighbor}
\end{equation}
kde \(U_t\) označuje množinu doposiaľ nenavštívených vrcholov.
Takto zostavený cyklus je výpočtovo lacný, no jeho kvalita závisí od
skorých lokálnych rozhodnutí \parencite{matai2010overview}. Preto sa
v druhom kroku uplatnilo lokálne zlepšovanie 2-opt
\parencite{croes1958twoopt}: pre každú dvojicu nepriľahlých hrán
\((a,b)\) a \((c,d)\) sa testuje podmienka
\begin{equation}
d_{ab} + d_{cd} > d_{ac} + d_{bd},
\label{eq:two_opt_condition}
\end{equation}
a ak je splnená, príslušný úsek trasy sa obráti. Postup sa opakuje,
kým v okolí aktuálneho riešenia existujú zlepšujúce výmeny.

\Needspace{8\baselineskip}
\indent V rovniciach~\eqref{eq:nearest_neighbor}
a~\eqref{eq:two_opt_condition} platí:
\begin{symboltable}
    \symbolitem{\(\pi^{(0)}_t\)}{označuje \(t\)-ty vrchol počiatočnej trasy pred zlepšovaním.}
    \symbolitem{\(U_t\)}{označuje množinu doposiaľ nenavštívených vrcholov v kroku \(t\).}
    \symbolitem{\(d_{ab},\, d_{cd}\)}{označujú pôvodné dĺžky dvoch vybraných hrán.}
    \symbolitem{\(d_{ac},\, d_{bd}\)}{označujú dĺžky hrán po ich vzájomnom nahradení.}
\end{symboltable}

\indent \textit{Presná referenčná úroveň.}
Popri heuristickej vetve sa vytvorila aj presná vetva, v ktorej sa
referenčné riešenia získavali riešičom Concorde
\parencite{concorde2003solver}. Keďže Concorde pracuje s diskrétnou
euklidovskou metrikou, súradnice sa pred exportom previedli na
celočíselnú reprezentáciu \(\tilde{x}_i = \lfloor s\, x_i \rceil\)
a vzdialenosti sa počítali ako
\begin{equation}
\tilde{d}_{ij}
= \left\lfloor \lVert \tilde{x}_i - \tilde{x}_j \rVert_2 + 0{,}5
  \right\rfloor,
\label{eq:tsplib_distance}
\end{equation}
čo zodpovedá diskrétnej metrike \texttt{EUC\_2D} používanej
v knižnici TSPLIB \parencite{reinelt1991tsplib}. Do presnej
vetvy sa zaradili iba prípady, pri ktorých Concorde potvrdil
optimálnosť nájdeného riešenia.

\Needspace{8\baselineskip}
\indent V tomto zápise platí:
\begin{symboltable}
    \symbolitem{\(x_i\)}{označuje pôvodnú normalizovanú súradnicu mesta.}
    \symbolitem{\(\tilde{x}_i\)}{označuje celočíselný obraz súradnice po zväčšení mierky a zaokrúhlení.}
    \symbolitem{\(s\)}{označuje škálovací faktor určujúci jemnosť celočíselnej mriežky.}
    \symbolitem{\(\lfloor \cdot \rceil\)}{označuje zaokrúhlenie na najbližšie celé číslo.}
    \symbolitem{\(\tilde{d}_{ij}\)}{označuje diskrétnu vzdialenosť medzi mestami \(i\) a \(j\).}
\end{symboltable}

\indent Po vygenerovaní sa pri každej inštancii ukladali súradnice,
referenčná trasa, jej dĺžka a sprievodné metadáta o pôvode riešenia.
Binárne štítky hrán sa nevytvárali v tejto fáze, ale až pri načítaní
inštancie do trénovacieho dátového radu — deterministicky z uloženej
permutácie.

%% ─── 3.3.2 TSPLIB ──────────────────────────────────────────────────────
\subsection{Štandardné testovacie inštancie TSPLIB}
\label{sec:tsplib}

\indent Záverečné hodnotenie modelov prebiehalo výhradne na štrnástich
štandardných inštanciách z knižnice TSPLIB
\parencite{reinelt1991tsplib}. Tieto inštancie neboli súčasťou
trénovacej ani validačnej množiny — slúžili ako nezávislý referenčný
rámec spoločný pre všetky porovnávané architektúry aj heuristické bázy,
čo umožnilo sledovať prenos naučených vzorov na úlohy odlišné od
synteticky generovaných príkladov.

\indent Súbor bol zostavený tak, aby pokrýval rôzne veľkosti
a geometrické štruktúry, a vyhol sa tak hodnoteniu naviazanému na
jeden úzky typ problému.

\begin{table}[H]
    \centering
    \small
    \caption{Štandardné testovacie inštancie použité pri záverečnom hodnotení.}
    \label{tab:tsplib_instances}
    \begin{tabular}{|p{3.0cm}|p{1.0cm}|p{3.0cm}|p{1.0cm}|}
        \hline
        Inštancia & \(n\) & Inštancia & \(n\) \\
        \hline
        \texttt{ulysses16} & 16 & \texttt{eil76}   & 76  \\
        \hline
        \texttt{ulysses22} & 22 & \texttt{pr76}    & 76  \\
        \hline
        \texttt{att48}     & 48 & \texttt{gr96}    & 96  \\
        \hline
        \texttt{eil51}     & 51 & \texttt{kroA100} & 100 \\
        \hline
        \texttt{berlin52}  & 52 & \texttt{kroC100} & 100 \\
        \hline
        \texttt{st70}      & 70 & \texttt{kroD100} & 100 \\
        \hline
        \texttt{rd100}     & 100 & \texttt{lin105} & 105 \\
        \hline
    \end{tabular}
    \normalsize
\end{table}

\indent Spracovanie každej inštancie začínalo načítaním súradníc
zo súboru \texttt{.tsp}. Keďže modely pracovali s normalizovaným
priestorom, pôvodné súradnice
\(x_i^{\mathrm{orig}} \in \mathbb{R}^2\) sa previedli afinnou
normalizáciou
\begin{equation}
x_i = \frac{x_i^{\mathrm{orig}} - m}{h},
\qquad
m = \min_{1 \leq k \leq n} x_k^{\mathrm{orig}},
\qquad
h = \max\!\left(\max_{1 \leq k \leq n} x_k^{\mathrm{orig}} - m,\;
\varepsilon\right),
\label{eq:tsplib_normalization}
\end{equation}
kde operácie \(\min\), \(\max\) a delenie sa vykonávali po
súradnicových osiach. Každá inštancia sa tým previedla do
normalizovaného obdĺžnika so stranami v intervale \([0,1]\).

\Needspace{8\baselineskip}
\indent V rovnici~\eqref{eq:tsplib_normalization} platí:
\begin{symboltable}
    \symbolitem{\(x_i^{\mathrm{orig}}\)}{označuje pôvodnú súradnicu mesta podľa definície inštancie TSPLIB.}
    \symbolitem{\(x_i\)}{označuje normalizovanú súradnicu použitú na vstupe do modelu.}
    \symbolitem{\(m\)}{označuje vektor minimálnych hodnôt pôvodných súradníc.}
    \symbolitem{\(h\)}{označuje vektor rozsahov po normalizácii, chránený konštantou \(\varepsilon\) proti nulovému deleniu.}
    \symbolitem{\(\varepsilon\)}{označuje malú kladnú konštantu použitú na numerickú stabilizáciu.}
\end{symboltable}

\indent Pôvodné súradnice sa pri spracovaní neodhadzovali — uchovávali
sa spolu s informáciou o metrike danej inštancie. Vďaka tomu bolo
možné pri záverečnom hodnotení počítať dĺžku trasy v pôvodnej
diskrétnej metrike TSPLIB. Referenčná trasa sa získavala prednostne
zo súboru \texttt{.opt.tour}; ak takýto súbor chýbal, použila sa
náhradná heuristická trasa.

\indent Ku každej spracovanej inštancii vznikol súbor \texttt{.npz}
obsahujúci normalizované aj pôvodné súradnice, referenčnú trasu, jej
dĺžku v oboch metrikách, typ metriky a parametre normalizácie.

%% ═══════════════════════════════════════════════════════════════════════════
%% 3.4  GRAFOVÁ REPREZENTÁCIA
%% ═══════════════════════════════════════════════════════════════════════════
\section{Grafová reprezentácia}
\label{sec:grafova_reprezentacia}

\indent Každá inštancia TSP s \(n\) mestami sa modelovala ako \emph{úplný
neorientovaný graf} \(G = (V, E)\), kde \(V = \{1, \ldots, n\}\) je množina
vrcholov (miest) a \(E = \{\{i,j\} \mid i < j\}\) je množina všetkých
\(m = n(n-1)/2\) hrán. Cieľom modelu nebolo vytvoriť trasu priamo, ale
priradiť každej hrane \emph{skóre} vyjadrujúce pravdepodobnosť jej
príslušnosti k optimálnemu Hamiltonovskému cyklu. Tento prístup prekladá
kombinatorický problém na problém binárnej klasifikácie nad hranami
a umožňuje vektorizované spracovanie celej inštancie v jednej dávke.

\indent Niektoré prístupy k dátovo riadeným metódam pre TSP pracujú
s riedkou grafovou reprezentáciou — zahŕňajú iba \(k\) najbližších susedov
každého vrcholu a tým znižujú počet hrán z \(\mathcal{O}(n^2)\) na
\(\mathcal{O}(kn)\) \parencite{joshi2022generalization}. V tejto
práci sa zachoval plný kandidátny priestor pre všetky porovnávané
architektúry aj heuristické bázy, pretože pre maximálnu veľkosť
inštancií (\(n = 105\)) je \(m_{\max} = 5\,460\), čo je počet hrán
zvládnuteľný v jednej dávke bez agresívneho orezávania.

%% ─── 3.4.1  Vektory vstupných príznakov hrán ────────────────────────────
\subsection{Vektory vstupných príznakov hrán}
\label{sec:hranove_priznaky}

\indent Pre každú neorientovanú hranu \(\{i,j\}\) s \(i < j\) sa zostavil
pevný desaťrozmerný vektor vstupných príznakov
\begin{equation}
\phi_{ij} =
\bigl[
  u_i,\;
  v_i,\;
  u_j,\;
  v_j,\;
  d_{ij},\;
  \Delta x_{ij},\;
  \Delta y_{ij},\;
  \alpha_{ij},\;
  (\Delta x_{ij})^2,\;
  (\Delta y_{ij})^2
\bigr]
\in \mathbb{R}^{10},
\label{eq:edge_features}
\end{equation}
kde jednotlivé zložky sú definované vzťahmi
\begin{align}
d_{ij} &= \lVert x_j - x_i \rVert_2,
\label{eq:feat_dist}\\
\Delta x_{ij} &= u_j - u_i, \qquad \Delta y_{ij} = v_j - v_i,
\label{eq:feat_delta}\\
\alpha_{ij} &= \operatorname{atan2}(\Delta y_{ij},\, \Delta x_{ij}).
\label{eq:feat_angle}
\end{align}

\begin{symboltable}
    \symbolitem{\((u_i, v_i)\)}{označuje normalizované súradnice vrcholu \(i\).}
    \symbolitem{\((u_j, v_j)\)}{označuje normalizované súradnice vrcholu \(j\).}
    \symbolitem{\(d_{ij}\)}{označuje euklidovskú vzdialenosť medzi vrcholmi \(i\) a \(j\).}
    \symbolitem{\(\Delta x_{ij},\, \Delta y_{ij}\)}{označujú horizontálnu a vertikálnu zmenu medzi vrcholmi.}
    \symbolitem{\(\alpha_{ij}\)}{označuje orientovaný uhol spojnice \((i,j)\) vzhľadom na os \(x\), hodnoty v \((-\pi, \pi]\).}
    \symbolitem{\((\Delta x_{ij})^2,\, (\Delta y_{ij})^2\)}{označujú kvadratické zložky zdôrazňujúce veľkosť zmien v oboch smeroch.}
    \symbolitem{\(\phi_{ij}\)}{označuje výsledný desaťrozmerný vektor príznakov hrany \(\{i,j\}\).}
\end{symboltable}

\indent Štruktúra vektora \(\phi_{ij}\) kombinuje dve skupiny informácií.
Prvé štyri zložky zachytávajú \emph{absolútnu polohu} oboch vrcholov,
stredná skupina popisuje \emph{relatívnu geometriu} spojenia a posledné
dve zložky zdôrazňujú kvadratickú škálu horizontálnych a vertikálnych
zmien. Všetky tri porovnávané rodiny modelov dostávali rovnaký vektor
\(\phi_{ij}\), čím sa zabezpečilo, že rozdiely vo výkone odrážajú
schopnosť architektúry spracovať informáciu, nie jej odlišný objem.

%% ─── 3.4.2  Odvodenie cieľových štítkov ─────────────────────────────────
\subsection{Odvodenie binárnych cieľových štítkov}
\label{sec:cielove_stitky}

\indent Posledným krokom prípravy vstupných dát bolo odvodenie binárneho
štítka pre každú hranu. Z referenčnej trasy \(\pi^*\) sa pre každú
neorientovanú hranu \(\{i,j\}\) s \(i < j\) určil štítok
\begin{equation}
y_{ij} =
\begin{cases}
1, & \text{ak } \{i,j\} \in E(\pi^*), \\
0, & \text{inak,}
\end{cases}
\label{eq:edge_label}
\end{equation}
kde \(E(\pi^*)\) označuje množinu hrán referenčného Hamiltonovského
cyklu. Tým sa úloha konštrukcie trasy previedla na binárnu klasifikáciu
nad hranami grafu: model sa neučil generovať poradie vrcholov, ale
odhadovať príslušnosť každého spojenia k výslednému cyklu.

\begin{symboltable}
    \symbolitem{\(\pi^*\)}{označuje referenčnú Hamiltonovskú trasu danej inštancie.}
    \symbolitem{\(E(\pi^*)\)}{označuje množinu neorientovaných hrán tvoriacich referenčný cyklus.}
    \symbolitem{\(y_{ij}\)}{označuje binárny cieľový štítok hrany \(\{i,j\}\).}
\end{symboltable}

\begin{figure}[H]
    \centering
    \resizebox{0.72\textwidth}{!}{% edge_label_derivation.tex
% K5 v pentagonálnom rozmiestnení. Tour = vonkajší päťuholník (y=1),
% diagonály = non-tour (y=0). Len legenda, bez inline štítkov.
% Použitie: \resizebox{0.55\textwidth}{!}{% edge_label_derivation.tex
% K5 v pentagonálnom rozmiestnení. Tour = vonkajší päťuholník (y=1),
% diagonály = non-tour (y=0). Len legenda, bez inline štítkov.
% Použitie: \resizebox{0.55\textwidth}{!}{\input{figures/tikz/edge_label_derivation}}
\begin{tikzpicture}[
    font=\small,
    vtx/.style={
        circle, draw=black!70, fill=white,
        minimum size=0.62cm, inner sep=1pt,
        font=\small\bfseries
    },
    touredge/.style={draw=blue!70, very thick},
    nonedge/.style={draw=gray!45, dashed, thin}
]

%% Vrcholy pravidelného päťuholníka r=2.1, stred=(2.60, 2.50)
\coordinate (v1) at (2.60, 4.60);
\coordinate (v2) at (4.60, 3.17);
\coordinate (v3) at (3.85, 0.90);
\coordinate (v4) at (1.35, 0.90);
\coordinate (v5) at (0.60, 3.17);

%% ── Diagonály (non-tour, y=0) ────────────────────────────────────
\draw[nonedge] (v1) -- (v3);
\draw[nonedge] (v1) -- (v4);
\draw[nonedge] (v2) -- (v4);
\draw[nonedge] (v2) -- (v5);
\draw[nonedge] (v3) -- (v5);

%% ── Vonkajší päťuholník (tour, y=1) ─────────────────────────────
\draw[touredge] (v1) -- (v2);
\draw[touredge] (v2) -- (v3);
\draw[touredge] (v3) -- (v4);
\draw[touredge] (v4) -- (v5);
\draw[touredge] (v5) -- (v1);

%% ── Vrcholy ──────────────────────────────────────────────────────
\node[vtx] at (v1) {1};
\node[vtx] at (v2) {2};
\node[vtx] at (v3) {3};
\node[vtx] at (v4) {4};
\node[vtx] at (v5) {5};

%% ── Legenda ──────────────────────────────────────────────────────
\draw[touredge] (5.20, 3.50) -- (6.00, 3.50);
\node[font=\scriptsize, anchor=west] at (6.10, 3.50)
    {hrana v trase (\(y_{ij}{=}1\))};
\draw[nonedge]  (5.20, 2.90) -- (6.00, 2.90);
\node[font=\scriptsize, anchor=west] at (6.10, 2.90)
    {nepatriaca hrana (\(y_{ij}{=}0\))};

\end{tikzpicture}}
\begin{tikzpicture}[
    font=\small,
    vtx/.style={
        circle, draw=black!70, fill=white,
        minimum size=0.62cm, inner sep=1pt,
        font=\small\bfseries
    },
    touredge/.style={draw=blue!70, very thick},
    nonedge/.style={draw=gray!45, dashed, thin}
]

%% Vrcholy pravidelného päťuholníka r=2.1, stred=(2.60, 2.50)
\coordinate (v1) at (2.60, 4.60);
\coordinate (v2) at (4.60, 3.17);
\coordinate (v3) at (3.85, 0.90);
\coordinate (v4) at (1.35, 0.90);
\coordinate (v5) at (0.60, 3.17);

%% ── Diagonály (non-tour, y=0) ────────────────────────────────────
\draw[nonedge] (v1) -- (v3);
\draw[nonedge] (v1) -- (v4);
\draw[nonedge] (v2) -- (v4);
\draw[nonedge] (v2) -- (v5);
\draw[nonedge] (v3) -- (v5);

%% ── Vonkajší päťuholník (tour, y=1) ─────────────────────────────
\draw[touredge] (v1) -- (v2);
\draw[touredge] (v2) -- (v3);
\draw[touredge] (v3) -- (v4);
\draw[touredge] (v4) -- (v5);
\draw[touredge] (v5) -- (v1);

%% ── Vrcholy ──────────────────────────────────────────────────────
\node[vtx] at (v1) {1};
\node[vtx] at (v2) {2};
\node[vtx] at (v3) {3};
\node[vtx] at (v4) {4};
\node[vtx] at (v5) {5};

%% ── Legenda ──────────────────────────────────────────────────────
\draw[touredge] (5.20, 3.50) -- (6.00, 3.50);
\node[font=\scriptsize, anchor=west] at (6.10, 3.50)
    {hrana v trase (\(y_{ij}{=}1\))};
\draw[nonedge]  (5.20, 2.90) -- (6.00, 2.90);
\node[font=\scriptsize, anchor=west] at (6.10, 2.90)
    {nepatriaca hrana (\(y_{ij}{=}0\))};

\end{tikzpicture}}
    \caption{Odvodenie binárnych štítkov z referenčnej trasy na príklade
             úplného grafu \(K_5\): hrany turu sú označené \(y_{ij}=1\),
             ostatné \(y_{ij}=0\).}
    \label{fig:edge_label_derivation}
\end{figure}

\indent Keďže Hamiltonovský cyklus obsahuje presne \(n\) hrán a úplný
graf ich má \(m = n(n-1)/2\), pozitívnych štítkov je výrazne menej ako
negatívnych. Tento nepomer priamo motivuje použitie váženej straty pri
trénovaní modelov, opísanej v sekcii~\ref{sec:trenovanie}.

%% ═══════════════════════════════════════════════════════════════════════════
%% 3.5  MODELOVÉ ARCHITEKTÚRY
%% ═══════════════════════════════════════════════════════════════════════════
\section{Modelové architektúry}
\label{sec:modelove_architektury}

\indent V tejto práci boli navrhnuté a porovnané tri rodiny neurónových
modelov. Všetky plnia rovnakú úlohu: na základe vektora geometrických
príznakov hrany \(\phi_{ij}\) odhadnúť skóre \(s_{ij}\), ktoré vyjadruje,
nakoľko pravdepodobne daná hrana patrí do výslednej trasy. Rodiny sa líšia
tým, \emph{koľko kontextu} o celej inštancii môže každý model pri tomto
odhade využiť.

\indent Neurónové siete sú trieda modelov strojového učenia inšpirovaná
spôsobom, akým biologické neuróny spracovávajú a prenášajú signály.
Základnou stavebnou jednotkou je \emph{umelý neurón} — funkcia, ktorá
prijíma niekoľko vstupných hodnôt, vypočíta ich váženú sumu a výsledok
transformuje nelineárnou \emph{aktivačnou funkciou}. Tá je kľúčová:
bez nej by sieť akejkoľvek hĺbky zostávala ekvivalentná jednej lineárnej
transformácii a nebola by schopná zachytiť zložité vzory v dátach.
V tejto práci sa používajú dve aktivačné funkcie: \textbf{ReLU},
ktorá záporné hodnoty nahradí nulou, a \textbf{GELU}, jej hladká
aproximácia s lepšími vlastnosťami pri trénovaní hlbokých sietí.

\indent Neuróny sa usporadúvajú do \emph{vrstiev}. Každá vrstva prijíma
výstupy predchádzajúcej vrstvy, transformuje ich a odovzdáva ďalej. Sieť
zložená z viacerých vrstiev za sebou sa nazýva \emph{viacvrstvový
perceptrón} (MLP). Čím viac vrstiev, tým abstraktnejšie vzory môže model
zachytiť — nižšie vrstvy sa učia jednoduché lokálne vlastnosti, vyššie
ich kombinujú do zložitejších konceptov.

\indent Váhy siete sa nastavujú počas \emph{trénovania}: model opakuje
predikcie nad trénovacími príkladmi, meria odchýlku od správnych hodnôt
pomocou \emph{stratovej funkcie} a algoritmom \emph{spätného šírenia}
posúva váhy smerom, ktorý chybu znižuje.
\textbf{Vrstvová normalizácia} (LayerNorm) priebežne upravuje rozsah
aktivácií v hlbokých sieťach a \textbf{Dropout} počas trénovania
náhodne deaktivuje časť neurónov, čím znižuje riziko preučenia.

\indent V kontexte tejto práce prijíma každý model vektor príznakov
hrany \(\phi_{ij}\) a vydáva jedno číslo — skóre \(s_{ij}\). Tri
porovnávané rodiny sa líšia predovšetkým tým, či pri hodnotení hrany
uvažujú len o nej samotnej, alebo zohľadňujú aj informáciu z celej
inštancie. Tabuľka~\ref{tab:arch_porovnanie} sumarizuje kľúčové
vlastnosti každej z nich.

\begin{table}[H]
\centering
\caption{Porovnanie troch porovnávaných architektúr neurónových modelov.}
\label{tab:arch_porovnanie}
\begin{tabular}{lp{2.6cm}p{2.6cm}p{3.2cm}}
\hline
\textbf{Vlastnosť}       & \textbf{EdgeMLP}               & \textbf{EdgeResMLP}            & \textbf{EdgeTransformer}         \\ \hline
Vstup                    & \(\phi_{ij}\in\mathbb{R}^{10}\) & \(\phi_{ij}\in\mathbb{R}^{10}\) & \(\phi_{ij}\in\mathbb{R}^{10}\)  \\
Spracovanie hrany        & lokálne, nezávislé             & lokálne, nezávislé             & globálne + lokálne               \\
Kontext inštancie        & žiadny                         & žiadny                         & celá inštancia                   \\
Aktivácia                & ReLU                           & GELU                           & GELU                             \\
Reziduálne spojenia      & nie                            & áno                            & áno (attn + FFN)                 \\
Hĺbka (vrstvy / bloky)  & 2 alebo 3                      & 4                              & 3                                \\
Skryté dimenzie          & 128 alebo 256                  & 128                            & 128                              \\
Hlavy pozornosti         & —                              & —                              & 4                                \\ \hline
\end{tabular}
\end{table}

%% ─── 3.5.1  Viacvrstvový perceptrón nad hranami ──────────────────────────
\subsection{Viacvrstvový perceptrón nad hranami}
\label{sec:edge_mlp}

\indent Najjednoduchšiu rodinu tvoril hranový MLP, v ktorom sa každá
hrana spracovávala úplne nezávisle od všetkých ostatných. Vektor
príznakov hrany \(\phi_{ij}\) postupne prešiel niekoľkými vrstvami
lineárnych transformácií s aktiváciou ReLU. Každá vrstva kombinuje
vstupné hodnoty, aplikuje nelinearitu a odovzdáva výsledok ďalšej.
Po poslednej skrytej vrstve záverečná lineárna projekcia vyprodukovala
skalárne skóre hrany \(s_{ij}\).

\begin{figure}[H]
    \centering
    \resizebox{0.72\textwidth}{!}{% mlp_edge_diagram.tex
% Schéma hranového MLP modelu (EdgeMLP)
% Použitie: \resizebox{0.72\textwidth}{!}{% mlp_edge_diagram.tex
% Schéma hranového MLP modelu (EdgeMLP)
% Použitie: \resizebox{0.72\textwidth}{!}{\input{figures/tikz/mlp_edge_diagram}}
\begin{tikzpicture}[
    font=\small,
    box/.style={
        draw, rounded corners=5pt, align=center,
        minimum width=3.4cm, minimum height=0.82cm,
        inner sep=4pt
    },
    inputbox/.style={box, fill=gray!10, draw=gray!50},
    layerbox/.style={box, fill=blue!8, draw=blue!40},
    outbox/.style={box, fill=green!8, draw=green!50},
    labelbox/.style={font=\scriptsize\itshape, text=gray!60, align=left,
                     text width=4.0cm},
    arrow/.style={-{Latex[length=2.5mm]}, thick}
]

\node[inputbox] (phi) at (0, 0)
    {\(\phi_{ij} \in \mathbb{R}^{10}\)\\[-2pt]
     \scriptsize vektor príznakov hrany};

\node[layerbox] (h1) at (0, -1.55)
    {Lineárna vrstva\\[-2pt]
     \scriptsize + aktivácia ReLU};

\node[layerbox] (hdots) at (0, -3.10)
    {\(\cdots\)\\[-2pt]
     \scriptsize (opakuje sa \(D\)-krát)};

\node[layerbox] (hD) at (0, -4.65)
    {Lineárna vrstva\\[-2pt]
     \scriptsize + aktivácia ReLU};

\node[outbox] (out) at (0, -6.15)
    {Záverečná projekcia\\[-2pt]
     \scriptsize logit \(s_{ij} \in \mathbb{R}\)};

\draw[arrow] (phi)   -- (h1);
\draw[arrow] (h1)    -- (hdots);
\draw[arrow] (hdots) -- (hD);
\draw[arrow] (hD)    -- (out);

\node[labelbox, anchor=north west] at (2.5, -0.25)
    {\textbf{Vstup:} 10-rozmerný vektor
     obsahujúci polohu oboch vrcholov,
     vzdialenosť, relatívnu geometriu
     a uhol hrany.};

\node[labelbox, anchor=north west] at (2.5, -2.00)
    {\textbf{Skryté vrstvy:} každá vrstva
     lineárne transformuje vstup
     a aplikuje ReLU nelinearitu.};

\node[labelbox, anchor=north west] at (2.5, -5.25)
    {\textbf{Výstup:} jediné číslo
     (logit) vyjadrujúce
     odhadovanú príslušnosť
     hrany k optimálnej trase.};

\node[draw=gray!35, rounded corners=4pt, fill=gray!5,
      font=\scriptsize, align=center, text width=6.5cm,
      inner sep=5pt] at (0, -7.55)
    {Každá hrana \(\{i,j\}\) sa spracováva \textbf{nezávisle} —
     model pri hodnotení hrany nemá prístup k informácii
     o ostatných hranách ani vrcholoch inštancie.};

\end{tikzpicture}
}
\begin{tikzpicture}[
    font=\small,
    box/.style={
        draw, rounded corners=5pt, align=center,
        minimum width=3.4cm, minimum height=0.82cm,
        inner sep=4pt
    },
    inputbox/.style={box, fill=gray!10, draw=gray!50},
    layerbox/.style={box, fill=blue!8, draw=blue!40},
    outbox/.style={box, fill=green!8, draw=green!50},
    labelbox/.style={font=\scriptsize\itshape, text=gray!60, align=left,
                     text width=4.0cm},
    arrow/.style={-{Latex[length=2.5mm]}, thick}
]

\node[inputbox] (phi) at (0, 0)
    {\(\phi_{ij} \in \mathbb{R}^{10}\)\\[-2pt]
     \scriptsize vektor príznakov hrany};

\node[layerbox] (h1) at (0, -1.55)
    {Lineárna vrstva\\[-2pt]
     \scriptsize + aktivácia ReLU};

\node[layerbox] (hdots) at (0, -3.10)
    {\(\cdots\)\\[-2pt]
     \scriptsize (opakuje sa \(D\)-krát)};

\node[layerbox] (hD) at (0, -4.65)
    {Lineárna vrstva\\[-2pt]
     \scriptsize + aktivácia ReLU};

\node[outbox] (out) at (0, -6.15)
    {Záverečná projekcia\\[-2pt]
     \scriptsize logit \(s_{ij} \in \mathbb{R}\)};

\draw[arrow] (phi)   -- (h1);
\draw[arrow] (h1)    -- (hdots);
\draw[arrow] (hdots) -- (hD);
\draw[arrow] (hD)    -- (out);

\node[labelbox, anchor=north west] at (2.5, -0.25)
    {\textbf{Vstup:} 10-rozmerný vektor
     obsahujúci polohu oboch vrcholov,
     vzdialenosť, relatívnu geometriu
     a uhol hrany.};

\node[labelbox, anchor=north west] at (2.5, -2.00)
    {\textbf{Skryté vrstvy:} každá vrstva
     lineárne transformuje vstup
     a aplikuje ReLU nelinearitu.};

\node[labelbox, anchor=north west] at (2.5, -5.25)
    {\textbf{Výstup:} jediné číslo
     (logit) vyjadrujúce
     odhadovanú príslušnosť
     hrany k optimálnej trase.};

\node[draw=gray!35, rounded corners=4pt, fill=gray!5,
      font=\scriptsize, align=center, text width=6.5cm,
      inner sep=5pt] at (0, -7.55)
    {Každá hrana \(\{i,j\}\) sa spracováva \textbf{nezávisle} —
     model pri hodnotení hrany nemá prístup k informácii
     o ostatných hranách ani vrcholoch inštancie.};

\end{tikzpicture}
}
    \caption{Schéma hranového MLP modelu. Vektor príznakov hrany
             \(\phi_{ij}\) prechádza sekvenciou lineárnych vrstiev
             s aktiváciou ReLU. Každá hrana sa spracováva nezávisle
             — model nemá informáciu o zvyšku grafu.}
    \label{fig:mlp_edge_diagram}
\end{figure}

\indent Výhodou tohto prístupu bola výpočtová jednoduchosť: celú
inštanciu bolo možné spracovať ako jednu dávku vektorov naraz, bez
akéhokoľvek sekvenčného algoritmu. Nevýhodou bolo, že model pri
hodnotení každej hrany nemal prístup k žiadnej informácii o zvyšku
grafu. V experimentoch sa testovali dve varianty tejto rodiny —
plytší model s dvoma skrytými vrstvami (\texttt{edge\_mlp},
\(d_h = 128\)) a hlbší model s tromi vrstvami a väčšou kapacitou
(\texttt{edge\_mlp\_deep}, \(d_h = 256\)).

%% ─── 3.5.2  Reziduálny MLP model ─────────────────────────────────────────
\subsection{Reziduálny MLP model}
\label{sec:edge_res_mlp}

\indent Druhú rodinu tvoril reziduálny MLP, ktorý rozšíril základný
prístup o dve dôležité vylepšenia. Namiesto jednoduchej sekvencie vrstiev
sieť pracovala s \emph{reziduálnymi blokmi}: každý blok sa neučil celé
zobrazenie vstupu na výstup, ale iba \emph{korekciu} voči tomu, čo
sa na vstup bloku priviedlo. Výstup bloku bol teda súčtom jeho vstupu
a naučenej korekcie — odtiaľ názov reziduálna (zvyšková) sieť
\parencite{he2016resnet}. Vnútri každého bloku sa navyše uplatnila
vrstvová normalizácia pred transformáciou a náhodné vypúšťanie za ňou.

\begin{figure}[H]
    \centering
    \resizebox{0.80\textwidth}{!}{% res_mlp_diagram.tex
% Diagram reziduálneho MLP modelu s vysvetlením LayerNorm a Dropout
% Použitie: \resizebox{0.80\textwidth}{!}{% res_mlp_diagram.tex
% Diagram reziduálneho MLP modelu s vysvetlením LayerNorm a Dropout
% Použitie: \resizebox{0.80\textwidth}{!}{\input{figures/tikz/res_mlp_diagram}}
\begin{tikzpicture}[
    font=\small,
    box/.style={
        draw, rounded corners=5pt, align=center,
        minimum width=3.4cm, minimum height=0.82cm,
        inner sep=4pt
    },
    inputbox/.style={box, fill=gray!10, draw=gray!50},
    projbox/.style={box, fill=blue!8, draw=blue!40},
    normbox/.style={box, fill=yellow!10, draw=yellow!60!black},
    dropbox/.style={box, fill=red!6, draw=red!30},
    outbox/.style={box, fill=green!8, draw=green!50},
    labelbox/.style={font=\scriptsize\itshape, text=gray!65, align=left,
                     text width=4.2cm},
    arrow/.style={-{Latex[length=2.5mm]}, thick},
    resarrow/.style={-{Latex[length=2.5mm]}, thick, draw=green!55!black}
]

\node[inputbox] (phi) at (0, 0)
    {\(\phi_{ij} \in \mathbb{R}^{10}\)\\[-2pt]
     \scriptsize vektor príznakov hrany};

\node[projbox] (proj) at (0, -1.55)
    {Vstupná projekcia\\[-2pt]
     \scriptsize + aktivácia GELU};

\node[normbox] (ln1) at (0, -3.10)
    {Vrstvová normalizácia\\[-2pt]
     \scriptsize (LayerNorm)};

\node[projbox] (fc1) at (0, -4.50)
    {Lineárna transformácia\\[-2pt]
     \scriptsize + aktivácia GELU};

\node[dropbox] (drop) at (0, -5.85)
    {Náhodné vypúšťanie\\[-2pt]
     \scriptsize (Dropout)};

\node[projbox] (fc2) at (0, -7.20)
    {Lineárna transformácia};

\node[draw=green!55!black, fill=green!5, circle,
      minimum size=0.65cm, font=\bfseries\small] (plus) at (0, -8.55) {+};

\draw[draw=green!40, dashed, rounded corners=6pt]
    (-2.3, -2.55) rectangle (2.3, -7.85);
\node[font=\scriptsize\itshape, text=green!60!black]
    at (0, -2.38) {reziduálny blok (opakuje sa \(D\)-krát)};

\draw[resarrow]
    (proj.east) -- ++(1.4, 0) -- ++(0, -7.0) -- (plus.east);
\node[font=\scriptsize, text=green!60!black, rotate=90, anchor=south]
    at (3.0, -5.2) {reziduálna skratka};

\node[outbox] (out) at (0, -9.95)
    {Záverečná lineárna vrstva\\[-2pt]
     \scriptsize logit \(s_{ij} \in \mathbb{R}\)};

\draw[arrow] (phi)  -- (proj);
\draw[arrow] (proj) -- (ln1);
\draw[arrow] (ln1)  -- (fc1);
\draw[arrow] (fc1)  -- (drop);
\draw[arrow] (drop) -- (fc2);
\draw[arrow] (fc2)  -- (plus);
\draw[arrow] (plus) -- (out);

\node[labelbox, anchor=north west] at (3.5, -3.25)
    {\textbf{LayerNorm:} normalizuje aktivácie
     v rámci jedného vzorku — stabilizuje tréning
     a urýchľuje konvergenciu.};

\node[labelbox, anchor=north west] at (3.5, -5.85)
    {\textbf{Dropout:} náhodne deaktivuje
     časť neurónov — zabraňuje preučeniu.};

\node[labelbox, anchor=north west] at (3.5, -8.55)
    {\textbf{Reziduálna skratka:}
     gradient prúdi priamo, bez
     útlmu cez hlboké bloky.};

\end{tikzpicture}
}
\begin{tikzpicture}[
    font=\small,
    box/.style={
        draw, rounded corners=5pt, align=center,
        minimum width=3.4cm, minimum height=0.82cm,
        inner sep=4pt
    },
    inputbox/.style={box, fill=gray!10, draw=gray!50},
    projbox/.style={box, fill=blue!8, draw=blue!40},
    normbox/.style={box, fill=yellow!10, draw=yellow!60!black},
    dropbox/.style={box, fill=red!6, draw=red!30},
    outbox/.style={box, fill=green!8, draw=green!50},
    labelbox/.style={font=\scriptsize\itshape, text=gray!65, align=left,
                     text width=4.2cm},
    arrow/.style={-{Latex[length=2.5mm]}, thick},
    resarrow/.style={-{Latex[length=2.5mm]}, thick, draw=green!55!black}
]

\node[inputbox] (phi) at (0, 0)
    {\(\phi_{ij} \in \mathbb{R}^{10}\)\\[-2pt]
     \scriptsize vektor príznakov hrany};

\node[projbox] (proj) at (0, -1.55)
    {Vstupná projekcia\\[-2pt]
     \scriptsize + aktivácia GELU};

\node[normbox] (ln1) at (0, -3.10)
    {Vrstvová normalizácia\\[-2pt]
     \scriptsize (LayerNorm)};

\node[projbox] (fc1) at (0, -4.50)
    {Lineárna transformácia\\[-2pt]
     \scriptsize + aktivácia GELU};

\node[dropbox] (drop) at (0, -5.85)
    {Náhodné vypúšťanie\\[-2pt]
     \scriptsize (Dropout)};

\node[projbox] (fc2) at (0, -7.20)
    {Lineárna transformácia};

\node[draw=green!55!black, fill=green!5, circle,
      minimum size=0.65cm, font=\bfseries\small] (plus) at (0, -8.55) {+};

\draw[draw=green!40, dashed, rounded corners=6pt]
    (-2.3, -2.55) rectangle (2.3, -7.85);
\node[font=\scriptsize\itshape, text=green!60!black]
    at (0, -2.38) {reziduálny blok (opakuje sa \(D\)-krát)};

\draw[resarrow]
    (proj.east) -- ++(1.4, 0) -- ++(0, -7.0) -- (plus.east);
\node[font=\scriptsize, text=green!60!black, rotate=90, anchor=south]
    at (3.0, -5.2) {reziduálna skratka};

\node[outbox] (out) at (0, -9.95)
    {Záverečná lineárna vrstva\\[-2pt]
     \scriptsize logit \(s_{ij} \in \mathbb{R}\)};

\draw[arrow] (phi)  -- (proj);
\draw[arrow] (proj) -- (ln1);
\draw[arrow] (ln1)  -- (fc1);
\draw[arrow] (fc1)  -- (drop);
\draw[arrow] (drop) -- (fc2);
\draw[arrow] (fc2)  -- (plus);
\draw[arrow] (plus) -- (out);

\node[labelbox, anchor=north west] at (3.5, -3.25)
    {\textbf{LayerNorm:} normalizuje aktivácie
     v rámci jedného vzorku — stabilizuje tréning
     a urýchľuje konvergenciu.};

\node[labelbox, anchor=north west] at (3.5, -5.85)
    {\textbf{Dropout:} náhodne deaktivuje
     časť neurónov — zabraňuje preučeniu.};

\node[labelbox, anchor=north west] at (3.5, -8.55)
    {\textbf{Reziduálna skratka:}
     gradient prúdi priamo, bez
     útlmu cez hlboké bloky.};

\end{tikzpicture}
}
    \caption{Schéma reziduálneho MLP modelu. Vstup prechádza vstupnou
             projekciou a následne sekvenciou reziduálnych blokov.
             Každý blok aplikuje vrstvovú normalizáciu, dve lineárne
             transformácie s aktiváciou GELU a náhodné vypúšťanie.
             Reziduálna skratka sčíta vstup bloku s jeho výstupom.}
    \label{fig:res_mlp_diagram}
\end{figure}

\indent Reziduálne zapojenie prináša dva praktické prínosy. Po prvé,
gradient môže prúdiť reziduálnymi skratkami priamo naprieč celou
hĺbkou siete bez toho, aby sa rozptýlil alebo explodoval — to umožňuje
stabilné trénovanie aj pri väčšom počte blokov. Po druhé, vrstvová
normalizácia udržiava aktivácie v zdravom rozsahu. Rovnako ako základný
MLP, aj tento model zostával \emph{lokálnym}: každá hrana sa
spracovávala nezávisle a model nemal priamy prístup ku kontextu zvyšku
grafu. Model bol konfigurovaný so štyrmi reziduálnymi blokmi
a skrytou dimenziou \(d_h = 128\).

%% ─── 3.5.3  Transformerový model s hranovým výstupom ────────────────────
\subsection{Transformerový model s hranovým výstupom}
\label{sec:edge_transformer}

\indent Tretiu a najzložitejšiu rodinu tvoril kontextový transformerový
model, ktorý ako jediný pracoval s \emph{globálnym kontextom} celej
inštancie. Namiesto toho, aby hodnotil každú hranu izolovane, model
najprv zostavil reprezentácie všetkých vrcholov s ohľadom na ich
vzájomné vzťahy — a až potom ohodnotil každú hranu.

\indent Model pracoval v dvoch fázach. V prvej fáze sa zo súradníc
každého mesta zostavila jeho vnútorná reprezentácia, ktorá sa následne
opakovane aktualizovala \emph{mechanizmom vlastnej pozornosti}
(self-attention). Tento mechanizmus umožňuje každému vrcholu
„zobrať do úvahy" všetky ostatné vrcholy inštancie a podľa toho
upraviť svoju vlastnú reprezentáciu, inšpirujúc sa transformerovou
architektúrou \parencite{vaswani2017attention}.
Výsledkom je, že každý vrchol nesie \emph{kontextovú reprezentáciu}
zakódovávajúcu nielen jeho vlastnú polohu, ale aj jeho vzťah
ku všetkým ostatným mestám. V druhej fáze sa pre každú hranu zlúčili
kontextové reprezentácie oboch jej krajných vrcholov spolu
s pôvodnými geometrickými príznakmi hrany. Toto zlúčenie vstúpilo
do výstupnej siete, ktorá vyprodukovala skóre hrany \(s_{ij}\).

\begin{figure}[H]
    \centering
    \resizebox{0.95\textwidth}{!}{% transformer_diagram.tex
% Diagram transformerového modelu — dve fázy, konceptuálny
% Použitie: \resizebox{0.88\textwidth}{!}{% transformer_diagram.tex
% Diagram transformerového modelu — dve fázy, konceptuálny
% Použitie: \resizebox{0.88\textwidth}{!}{\input{figures/tikz/transformer_diagram}}
\begin{tikzpicture}[
    font=\small,
    box/.style={
        draw, rounded corners=5pt, align=center,
        minimum width=3.6cm, minimum height=0.85cm,
        inner sep=5pt
    },
    inputbox/.style={box, fill=gray!10, draw=gray!50},
    attnbox/.style={box, fill=orange!10, draw=orange!50},
    mergebox/.style={box, fill=purple!8, draw=purple!40},
    outbox/.style={box, fill=green!8, draw=green!50},
    labelbox/.style={font=\scriptsize\itshape, text=gray!65, align=left,
                     text width=5.0cm},
    arrow/.style={-{Latex[length=2.5mm]}, thick},
    dasharrow/.style={-{Latex[length=2mm]}, dashed, gray!60}
]

%% ══════════════════════════════════════════════════════════════════
%% FÁZA 1 — Kódovanie vrcholov
%% ══════════════════════════════════════════════════════════════════
\node[font=\bfseries\small, text=orange!70!black]
    at (0, 0.55) {Fáza 1 — kódovanie vrcholov};

\node[inputbox] (coords) at (0, -0.45)
    {Súradnice miest \((u_i, v_i)\)\\[-2pt]
     \scriptsize pre všetky vrcholy inštancie};

\node[attnbox] (embed) at (0, -2.00)
    {Lineárna projekcia\\[-2pt]
     \scriptsize súradnice → vnútorná reprezentácia};

\node[attnbox] (attn) at (0, -3.55)
    {Mechanizmus vlastnej pozornosti\\[-2pt]
     \scriptsize každý vrchol "sleduje" všetky ostatné};

\node[attnbox] (repeat) at (0, -5.10)
    {\(\cdots\) opakuje sa \(L\)-krát};

\node[outbox] (ctx) at (0, -6.55)
    {Kontextové reprezentácie \(z_i\)\\[-2pt]
     \scriptsize vrchol = poloha + vzťah k ostatným};

\draw[arrow] (coords) -- (embed);
\draw[arrow] (embed)  -- (attn);
\draw[arrow] (attn)   -- (repeat);
\draw[arrow] (repeat) -- (ctx);

%% Ohraničenie fázy 1
\draw[draw=orange!35, rounded corners=6pt, dashed]
    (-2.5, 0.20) rectangle (2.5, -7.20);

%% ══════════════════════════════════════════════════════════════════
%% FÁZA 2 — Skórovanie hrán
%% ══════════════════════════════════════════════════════════════════
\node[font=\bfseries\small, text=purple!70!black]
    at (7.0, 0.55) {Fáza 2 — skórovanie hrán};

%% Vstupné zdroje pre hranu {i,j}
\node[outbox] (zi) at (5.5, -0.45)
    {\(z_i\) — kontext vrcholu \(i\)};
\node[outbox] (zj) at (8.5, -0.45)
    {\(z_j\) — kontext vrcholu \(j\)};
\node[inputbox] (phi) at (7.0, -2.00)
    {\(\phi_{ij}\) — príznaky hrany};

%% Spojenie
\node[mergebox] (merge) at (7.0, -3.55)
    {Zlúčenie informácií\\[-2pt]
     \scriptsize kontext + príznaky hrany};

\node[mergebox] (head) at (7.0, -5.10)
    {Výstupná sieť\\[-2pt]
     \scriptsize (dve lineárne vrstvy + GELU)};

\node[outbox] (score) at (7.0, -6.55)
    {Skóre hrany \(s_{ij} \in \mathbb{R}\)\\[-2pt]
     \scriptsize miera príslušnosti k trase};

\draw[arrow] (zi)    |- (merge);
\draw[arrow] (zj)    |- (merge);
\draw[arrow] (phi)   -- (merge);
\draw[arrow] (merge) -- (head);
\draw[arrow] (head)  -- (score);

%% Prepojenie fázy 1 → fáza 2
\draw[arrow, draw=orange!60] (ctx.east) -- ++(1.0,0)
    |- node[midway, above, font=\scriptsize, text=orange!70]
       {pre každú hranu} (zi.west);
\draw[arrow, draw=orange!60] (ctx.east) -- ++(1.0,0)
    |- (zj.west);

%% Ohraničenie fázy 2
\draw[draw=purple!30, rounded corners=6pt, dashed]
    (4.2, 0.20) rectangle (9.8, -7.20);

%% ── Vysvetlivka dole ─────────────────────────────────────────────
\node[font=\scriptsize, text=gray!65, align=center, text width=13cm]
    at (3.5, -8.10)
    {Kľúčový rozdiel oproti MLP: transformer najprv zostaví
     \emph{globálny kontext} celej inštancie (Fáza 1) a až potom
     hodnotí každú hranu — s vedomosťou o jej vzťahu ku všetkým
     ostatným vrcholom (Fáza 2).};

\end{tikzpicture}}
\begin{tikzpicture}[
    font=\small,
    box/.style={
        draw, rounded corners=5pt, align=center,
        minimum width=3.6cm, minimum height=0.85cm,
        inner sep=5pt
    },
    inputbox/.style={box, fill=gray!10, draw=gray!50},
    attnbox/.style={box, fill=orange!10, draw=orange!50},
    mergebox/.style={box, fill=purple!8, draw=purple!40},
    outbox/.style={box, fill=green!8, draw=green!50},
    labelbox/.style={font=\scriptsize\itshape, text=gray!65, align=left,
                     text width=5.0cm},
    arrow/.style={-{Latex[length=2.5mm]}, thick},
    dasharrow/.style={-{Latex[length=2mm]}, dashed, gray!60}
]

%% ══════════════════════════════════════════════════════════════════
%% FÁZA 1 — Kódovanie vrcholov
%% ══════════════════════════════════════════════════════════════════
\node[font=\bfseries\small, text=orange!70!black]
    at (0, 0.55) {Fáza 1 — kódovanie vrcholov};

\node[inputbox] (coords) at (0, -0.45)
    {Súradnice miest \((u_i, v_i)\)\\[-2pt]
     \scriptsize pre všetky vrcholy inštancie};

\node[attnbox] (embed) at (0, -2.00)
    {Lineárna projekcia\\[-2pt]
     \scriptsize súradnice → vnútorná reprezentácia};

\node[attnbox] (attn) at (0, -3.55)
    {Mechanizmus vlastnej pozornosti\\[-2pt]
     \scriptsize každý vrchol "sleduje" všetky ostatné};

\node[attnbox] (repeat) at (0, -5.10)
    {\(\cdots\) opakuje sa \(L\)-krát};

\node[outbox] (ctx) at (0, -6.55)
    {Kontextové reprezentácie \(z_i\)\\[-2pt]
     \scriptsize vrchol = poloha + vzťah k ostatným};

\draw[arrow] (coords) -- (embed);
\draw[arrow] (embed)  -- (attn);
\draw[arrow] (attn)   -- (repeat);
\draw[arrow] (repeat) -- (ctx);

%% Ohraničenie fázy 1
\draw[draw=orange!35, rounded corners=6pt, dashed]
    (-2.5, 0.20) rectangle (2.5, -7.20);

%% ══════════════════════════════════════════════════════════════════
%% FÁZA 2 — Skórovanie hrán
%% ══════════════════════════════════════════════════════════════════
\node[font=\bfseries\small, text=purple!70!black]
    at (7.0, 0.55) {Fáza 2 — skórovanie hrán};

%% Vstupné zdroje pre hranu {i,j}
\node[outbox] (zi) at (5.5, -0.45)
    {\(z_i\) — kontext vrcholu \(i\)};
\node[outbox] (zj) at (8.5, -0.45)
    {\(z_j\) — kontext vrcholu \(j\)};
\node[inputbox] (phi) at (7.0, -2.00)
    {\(\phi_{ij}\) — príznaky hrany};

%% Spojenie
\node[mergebox] (merge) at (7.0, -3.55)
    {Zlúčenie informácií\\[-2pt]
     \scriptsize kontext + príznaky hrany};

\node[mergebox] (head) at (7.0, -5.10)
    {Výstupná sieť\\[-2pt]
     \scriptsize (dve lineárne vrstvy + GELU)};

\node[outbox] (score) at (7.0, -6.55)
    {Skóre hrany \(s_{ij} \in \mathbb{R}\)\\[-2pt]
     \scriptsize miera príslušnosti k trase};

\draw[arrow] (zi)    |- (merge);
\draw[arrow] (zj)    |- (merge);
\draw[arrow] (phi)   -- (merge);
\draw[arrow] (merge) -- (head);
\draw[arrow] (head)  -- (score);

%% Prepojenie fázy 1 → fáza 2
\draw[arrow, draw=orange!60] (ctx.east) -- ++(1.0,0)
    |- node[midway, above, font=\scriptsize, text=orange!70]
       {pre každú hranu} (zi.west);
\draw[arrow, draw=orange!60] (ctx.east) -- ++(1.0,0)
    |- (zj.west);

%% Ohraničenie fázy 2
\draw[draw=purple!30, rounded corners=6pt, dashed]
    (4.2, 0.20) rectangle (9.8, -7.20);

%% ── Vysvetlivka dole ─────────────────────────────────────────────
\node[font=\scriptsize, text=gray!65, align=center, text width=13cm]
    at (3.5, -8.10)
    {Kľúčový rozdiel oproti MLP: transformer najprv zostaví
     \emph{globálny kontext} celej inštancie (Fáza 1) a až potom
     hodnotí každú hranu — s vedomosťou o jej vzťahu ku všetkým
     ostatným vrcholom (Fáza 2).};

\end{tikzpicture}}
    \caption{Schéma transformerového modelu s hranovým výstupom.
             Fáza~1 zostrojí kontextové reprezentácie vrcholov
             pomocou mechanizmu vlastnej pozornosti.
             Fáza~2 zlúči tieto reprezentácie s hranovými príznakmi
             a výstupnou sieťou vyprodukuje skóre hrany.}
    \label{fig:transformer_diagram}
\end{figure}

\indent Kľúčovou výhodou tohto prístupu je, že pri rozhodovaní o každej
hrane model zohľadňuje celú štruktúru inštancie — nie len samotnú hranu.
Nevýhodou je vyššia výpočtová náročnosť: mechanizmus vlastnej pozornosti
porovnáva každý vrchol s každým iným (\(\mathcal{O}(n^2)\)), čo si
vyžaduje viac výpočtov ako jednoduché lokálne spracovanie hrán.
Model bol konfigurovaný s tromi transformerovými vrstvami,
\(d_h = 128\) a štyrmi hlavami pozornosti.

%% ═══════════════════════════════════════════════════════════════════════════
%% 3.6  TRÉNOVANIE
%% ═══════════════════════════════════════════════════════════════════════════
\section{Trénovanie}
\label{sec:trenovanie}

\indent Trénovací postup bol navrhnutý tak, aby bol reprodukovateľný
a konzistentný naprieč všetkými porovnávanými konfiguráciami. Každý model
sa trénoval z náhodne inicializovaných váh pomocou optimalizátora
Adam s počiatočnou mierou učenia \(10^{-3}\). Tréning prebiehal po
epochách, pričom po každej epoche sa validačná strata vyhodnotila na
vyhradenú validačnú množinu.

\indent Výber stratovej funkcie bol priamo motivovaný štruktúrou problému.
Keďže Hamiltonovský cyklus obsahuje presne \(n\) hrán z \(m = n(n-1)/2\)
celkových, je počet pozitívnych štítkov rádovo menší ako počet negatívnych
— tento nepomer narastá kvadraticky s veľkosťou inštancie. Bežná binárna
krížová entropia by preto model nevyhnutne tlačila k predpovedaniu samých
negatívnych štítkov. Z toho dôvodu bola použitá \emph{vážená binárna
krížová entropia} s triedami dynamicky vyváženými podľa skutočného pomeru
pozitívnych a negatívnych hrán v každej trénovacej dávke:

\begin{equation}
\mathcal{L} = -\frac{1}{|E|}
\sum_{\{i,j\} \in E}
\Bigl[
    w_+ \cdot y_{ij} \log \sigma(s_{ij})
    +
    w_- \cdot (1 - y_{ij}) \log \bigl(1 - \sigma(s_{ij})\bigr)
\Bigr],
\label{eq:bce_balanced}
\end{equation}

\noindent kde \(\sigma(\cdot)\) označuje sigmoidu, \(s_{ij}\) je surový
logit modelu a váhy \(w_+\), \(w_-\) sa vypočítajú z každej dávky ako

\begin{equation}
w_+ = \frac{|E^-|}{|E|}, \qquad w_- = \frac{|E^+|}{|E|},
\label{eq:bce_weights}
\end{equation}

\noindent pričom \(|E^+|\) a \(|E^-|\) označujú počet pozitívnych,
resp. negatívnych hrán v dávke a \(|E| = |E^+| + |E^-|\).

\begin{symboltable}
    \symbolitem{\(\mathcal{L}\)}{označuje hodnotu váženej binárnej krížovej entropie.}
    \symbolitem{\(|E|\)}{označuje celkový počet hrán v trénovacej dávke.}
    \symbolitem{\(w_+,\, w_-\)}{označujú váhy pozitívnej, resp. negatívnej triedy.}
    \symbolitem{\(\sigma(s_{ij})\)}{označuje sigmoidálnu transformáciu logitu hrany \(\{i,j\}\).}
\end{symboltable}

\indent Miera učenia sa znižovala automaticky pomocou plánovača
\textit{ReduceLROnPlateau}: ak sa validačná strata nezlepšila po dobu
desiatich po sebe nasledujúcich overovacích kôl, miera učenia sa
znásobila faktorom \(0{,}5\). Minimálna prípustná miera učenia bola
nastavená na \(10^{-5}\). Predčasné zastavenie nastalo v prípade,
že sa validačná strata nezlepšila ani po pätnástich overovacích kolách;
v tom momente sa tréning ukončil a uložil sa najlepší kontrolný bod.
Maximálny počet epoch bol nastavený na 80.

%% ═══════════════════════════════════════════════════════════════════════════
%% 3.7  DEKÓDOVANIE
%% ═══════════════════════════════════════════════════════════════════════════
\section{Dekódovanie}
\label{sec:dekodovanie}

\indent Neurónový model po natrénovaní vyprodukuje pre každú hranu
\(\{i,j\}\) surový logit \(s_{ij}\), no samotný logit nie je
priamo použiteľnou trasou — je potrebný krok \emph{dekódovania},
ktorý zo skóre hrán skonštruuje realizovateľný Hamiltonovský cyklus.

\indent Greedy cyklový dekodér postupuje nasledovne: hrany sa zoradia
zostupne podľa \(s_{ij}\) a prechádzajú sa v tomto poradí. Hrana
\(\{i,j\}\) sa zaradí do trasy, ak sú splnené tri podmienky: (1) stupeň
vrcholu \(i\) ani \(j\) ešte nedosiahol hodnotu 2; (2) zaradenie hrany
nevytvorí predčasný cyklus skôr, ako je v trase \(n\) hrán; (3) výsledná
podštruktúra ostáva Hamiltonovskou cestou. Podmienka (2) sa kontroluje
pomocou dátovej štruktúry \textit{union-find}. Po zaradení \(n\) hrán
vznikne Hamiltonovský cyklus prechádzajúci každým vrcholom práve raz.

\indent Na zlepšenie kvality získaných trás bola aplikovaná lokálna
optimalizácia \textbf{2-opt} \parencite{croes1958twoopt}. Táto procedúra
iteratívne prehľadáva všetky dvojice hrán a overuje, či ich výmena
skráti celkovú dĺžku trasy. Postup sa opakuje, kým sa nenájde žiadne
zlepšenie alebo kým sa nedosiahne nastavený limit počtu priebehov.

\indent Na posilnenie robustnosti dekódovania boli pre časť experimentov
použité dva dodatočné mechanizmy. \textbf{Multistart dekódovanie}
spúšťa greedy dekodér viackrát, pričom logity sa pred každým spustením
mierne narušia prídavkom gaussovského šumu so smerodajnou odchýlkou
\(\sigma_\text{noise}\). Z viacerých získaných trás sa vyberie
najkratšia. Konfigurácia so základným profilom používala
\(\text{multistart} = 1\) bez šumu, kým optimalizovaný profil
nastavoval \(\text{multistart} = 8\) a \(\sigma_\text{noise} = 0{,}02\).

%% ═══════════════════════════════════════════════════════════════════════════
%% 3.8  HODNOTENIE
%% ═══════════════════════════════════════════════════════════════════════════
\section{Hodnotenie}
\label{sec:hodnotenie}

\subsection{Metriky}
\label{sec:metriky}

\indent Primárnou metrikou hodnotenia kvality trás bola \emph{relatívna
odchýlka od optima} (optimality gap), vyjadrená v percentách:

\begin{equation}
\Delta(\%) = \frac{L_\text{model} - L_\text{ref}}{L_\text{ref}} \times 100,
\label{eq:gap}
\end{equation}

\noindent kde \(L_\text{model}\) je dĺžka trasy získanej modelom
a \(L_\text{ref}\) je dĺžka referenčnej trasy pre danú inštanciu.
Pre inštancie TSPLIB zodpovedá referenčná trasa optimálnemu riešeniu
uvedenému v knižnici \parencite{reinelt1991tsplib}.

\begin{symboltable}
    \symbolitem{\(\Delta(\%)\)}{označuje relatívnu odchýlku trasy modelu od referenčného optima.}
    \symbolitem{\(L_\text{model}\)}{označuje dĺžku trasy nájdenej modelom v pôvodnej metrike.}
    \symbolitem{\(L_\text{ref}\)}{označuje dĺžku referenčnej optimálnej trasy.}
\end{symboltable}

\indent Táto metrika je nezávislá od veľkosti inštancie a umožňuje
porovnanie naprieč rôznymi \(n\), čo bolo kľúčové pri vyhodnotení
18 experimentálnych konfigurácií na 14 inštanciách TSPLIB. Okrem
strednej hodnoty odchýlky sa sledovalo aj rozdelenie odchýlok naprieč
inštanciami, čo umožnilo identifikovať prípady, pri ktorých modely
systematicky zlyhávajú.

\subsection{Kontrola kvality a analytika}
\label{sec:qa_analytika}

\indent Pred spustením akéhokoľvek tréningu prebehla automatická
kontrola kvality (QA) nad všetkými vygenerovanými dátovými súbormi.
Skript overoval: (1) platnosť Hamiltonovského cyklu pre každú
referenčnú trasu; (2) zhodu dĺžky uloženého cyklu s dĺžkou
prepočítanou zo súradníc; (3) neprítomnosť poškodených záznamov.
Výsledky QA sú uložené v súboroch \texttt{runs/qa\_report.csv}
a \texttt{runs/qa\_report\_concorde.csv}; oba reporty zachytávajú
nula neplatných záznamov a nula neshôd v dĺžke.

\indent Po dokončení všetkých experimentálnych behov sa spustila
analytická fáza konsolidujúca výsledky zo 18 konfigurácií do súhrnných
tabuliek. Hlavnými výstupmi boli: celkový rebríček modelov podľa
strednej odchýlky (\texttt{model\_ranking.csv}); porovnanie základného
a optimalizovaného dekódovacieho profilu
(\texttt{eval\_profile\_compare.csv}); matica odchýlok na úrovni
jednotlivých inštancií TSPLIB (\texttt{instance\_gap\_matrix.csv})
a súhrnný prehľad priebehu tréningu (\texttt{experiments\_summary.csv}).
\chapter{Výsledky práce a diskusia}

\section{Prezentácia výsledkov}
\indent \textit{Sem doplňte text.}

\section{Diskusia výsledkov}
\indent \textit{Sem doplňte text.}

\section{Prínos práce}
\indent \textit{Sem doplňte text.}

\chapter{Záver}

\section{Zhrnutie dosiahnutých výsledkov}
\indent \textit{Sem doplňte text.}

\section{Odporúčania pre ďalší výskum}
\indent \textit{Sem doplňte text.}


\printbibliography[title={Zoznam použitej literatúry}]

\appendix
\chapter{Prílohy}

\noindent
\textit{Sem vložte prílohy práce (ak sú súčasťou práce).}


\end{document}
